\chapter{Temporal Trajectories and Segal Execution Semantics}
\label{ch:temporal-trajectories-segal-execution}

\section{Orientation: from transitions to finite executions}
\label{sec:temporal-from-mealy-executions}

The preceding chapters developed the state semantics of categorical automata.
A state of a Mealy morphism determines a residual behaviour, behavioural
equivalence identifies states with the same residual behaviour, and semantic
factorization produces canonical separated realizations. Those constructions
answer a state-oriented question:
\[
    \text{what behaviour is represented by the present state?}
\]

This chapter begins from the complementary execution-oriented question:
\[
    \text{what finite run is forced by a state and a finite input path?}
\]

For a Mealy morphism
\[
    \Phi=(P,\delta_P,\omega_P):
    \mathbb A\rightsquigarrow\mathbb B,
\]
the transition operation sends an applicable input arrow and a current state
to the next state, while the output operation assigns the corresponding
output-comparison arrow. Repeating the one-step operation produces a finite
labelled execution
\[
\begin{tikzcd}[column sep=large]
    x_0
        \arrow[r, "{a_1}"]
    &
    x_1
        \arrow[r, "{a_2}"]
    &
    \cdots
        \arrow[r]
    &
    x_n,
\end{tikzcd}
\]
together with an input path in \(\mathbb A^{\mathrm{op}}\) and an
output-comparison path in \(\mathbb B^{\mathrm{op}}\).

The constructions are organized into three successive presentations:
\[
\begin{tikzcd}[column sep=huge]
    \text{transition-output data}
        \arrow[r]
    &
    \operatorname{Exec}(\Phi)
        \arrow[r]
    &
    \mathscr T_\Phi.
\end{tikzcd}
\]
The execution category \(\operatorname{Exec}(\Phi)\) is the operational
presentation.  Its path object \(\mathscr T_\Phi\) is the finite-trajectory
presentation.  The following chapters will evaluate output-comparison labels
in downstream actions, transpose the resulting labelled actions into
polynomial total-response coalgebras, and then develop span, profile, and
trajectory predicate transformers.  The present chapter supplies the
execution geometry on which those later constructions depend.

Finite executions already exhibit the structures isolated below:
\begin{enumerate}
    \item a run restricts to every contiguous temporal window;
    \item adjacent runs glue uniquely when their common boundary states agree;
    \item an identity execution may be inserted at any boundary;
    \item adjacent execution steps may be contracted using composition in the
    input category;
    \item the Mealy multiplication laws ensure that state transition and
    output comparison are compatible with that contraction.
\end{enumerate}

The restriction and gluing operations belong to the path-sheaf structure.
Identity insertion and contraction use the internal category operations and
form the additional Segal execution structure.  This distinction will later
reappear, after dependent-product transposition, as the distinction between
an arbitrary one-step response coalgebra and a multiplicative response
coalgebra.

The principal results are organized in the order in which these structures
are derived:
\begin{enumerate}
    \item every Mealy morphism determines an execution category labelled over
    \(\mathbb A^{\mathrm{op}}\) and \(\mathbb B^{\mathrm{op}}\);
    \item finite executions are uniquely determined by an initial state and a
    finite input path;
    \item execution categories compose by pullback over their common labelled
    interface, naturally and coherently;
    \item the finite path objects of any internal graph form a sheaf on the
    discrete interval site, and internal graphs are equivalent to trajectory
    sheaves;
    \item internal category structures are exactly Segal composition
    structures on the corresponding trajectory sheaves;
    \item strict regular semantic quotients induce regular epimorphisms in
    every finite path length, so behavioural minimization induces
    finite-trajectory minimization;
    \item in canonical residual execution, derivative generation is finite-path
    endpoint reachability;
    \item in the Boolean specialization, the full word-labelled execution
    category is the free categorical completion of the letterwise derivative
    graph.
\end{enumerate}

The ambient hypotheses are used locally.  The execution category, finite-path
objects, interval site, trajectory-sheaf equivalence, and Segal constructions
require only finite limits.  The quotient and minimization results use the
regular or exact structure stated in their sections.  The canonical residual
application additionally uses the coproducts and regular images required to
form finite-path endpoint reachability.  Universal trajectory boxes and
safety cores require further dependent products and complete subobject
lattices; those logical constructions are deferred to the later dynamic-logic
chapter.

Unless a section states otherwise, \(\mathcal E\) is finitely complete and
chosen pullbacks are used throughout.  All external indexing categories lie
in a fixed universe.  Internal categories, Mealy morphisms, strict semantic
homomorphisms, semantic images, derivative closure, and the canonical
behaviour machine use the conventions established in the preceding chapters.

Write
\[
    \mathbb A=(A_0,A_1,s_A,t_A,e_A,m_A),
    \qquad
    \mathbb B=(B_0,B_1,s_B,t_B,e_B,m_B).
\]
The carrier of \(\Phi\) has anchors
\[
    p_P:P\to A_0,
    \qquad
    q_P:P\to B_0.
\]
The transition domain is
\[
    D_P
    :=
    A_1\times_{t_A,A_0,p_P}P,
\]
and the structure maps are written
\[
    \delta_P:D_P\to P,
    \qquad
    \omega_P:D_P\to B_1.
\]
For an arrow
\[
    a:u\to v
\]
and a state \(x\) with \(p_P(x)=v\), the transition and output comparison
have types
\[
    \delta_P(a,x)\in P_u,
\]
\[
    \omega_P(a,x):
    q_P\bigl(\delta_P(a,x)\bigr)
    \longrightarrow
    q_P(x).
\]
For composable arrows
\[
    a:u\to v,
    \qquad
    b:v\to w,
\]
and a state \(x\) over \(w\), the Mealy multiplication laws are
\[
    \delta_P(b\circ a,x)
    =
    \delta_P\bigl(a,\delta_P(b,x)\bigr),
\]
\[
    \omega_P(b\circ a,x)
    =
    \omega_P(b,x)
    \circ
    \omega_P\bigl(a,\delta_P(b,x)\bigr).
\]
These equations determine the orientation of every execution construction
below.

\section{The execution category of a Mealy morphism}
\label{subsec:mealy-action-category}

A one-step execution is not merely an input arrow. It consists of an input
arrow together with the state at which that arrow is processed. The next state
is then determined by the transition operation.

\begin{definition}[Execution category]
\label{def:execution-category}
The \textbf{execution category} of \(\Phi\) is the internal category
\[
    \operatorname{Exec}(\Phi)
\]
whose object object is
\[
    \operatorname{Exec}(\Phi)_0:=P
\]
and whose arrow object is
\[
    \operatorname{Exec}(\Phi)_1:=D_P.
\]
Its source and target maps are
\[
    s_{\operatorname{Exec}}(a,x):=x,
    \qquad
    t_{\operatorname{Exec}}(a,x):=\delta_P(a,x).
\]
\end{definition}

Thus a one-step execution has the form
\[
\begin{tikzcd}[column sep=huge]
    x
        \arrow[r, "{(a,x)}"]
    &
    \delta_P(a,x).
\end{tikzcd}
\]
The identity execution at \(x\) is
\[
    \bigl(e_A(p_Px),x\bigr).
\]
If
\[
    a:u\to v,
    \qquad
    b:v\to w,
\]
and \(x\) lies over \(w\), then the two execution arrows
\[
\begin{tikzcd}[column sep=large]
    x
        \arrow[r, "{(b,x)}"]
    &
    \delta_P(b,x)
        \arrow[r, "{(a,\delta_P(b,x))}"]
    &
    \delta_P\bigl(a,\delta_P(b,x)\bigr)
\end{tikzcd}
\]
are composable. Their composite is defined to be
\[
\begin{tikzcd}[column sep=huge]
    x
        \arrow[r, "{(b\circ a,x)}"]
    &
    \delta_P(b\circ a,x).
\end{tikzcd}
\]
The transition multiplication law identifies the two displayed terminal
states.

\begin{proposition}[Internal category structure of execution]
\label{prop:internal-category-structure-execution}
The displayed source, target, identity, and composition maps make
\(\operatorname{Exec}(\Phi)\) an internal category.
\end{proposition}

\begin{proof}
The source and target typing equations follow from the anchor equations of the
Mealy transition. The transition unit law gives
\[
    \delta_P(e_A(p_Px),x)=x,
\]
so the displayed execution arrow is an identity at \(x\).

For a composable pair, the transition multiplication equation
\[
    \delta_P(b\circ a,x)
    =
    \delta_P\bigl(a,\delta_P(b,x)\bigr)
\]
shows that the proposed composite has the same target as the successive
execution of \(b\) and then \(a\). The two unit laws for the execution
category are inherited from the unit laws of \(\mathbb A\) together with the
Mealy transition unit law.

For a composable triple of input arrows, both parenthesized execution
composites have input label equal to the composite arrow in \(\mathbb A\),
and their terminal states agree by repeated application of the transition
multiplication law. Associativity in \(\mathbb A\) therefore gives
associativity in \(\operatorname{Exec}(\Phi)\).
\end{proof}

The execution category is the action category of the contravariant input
action. With that notation one may write
\[
    \operatorname{Exec}(\Phi)
    \cong
    P\rtimes\mathbb A^{\mathrm{op}}.
\]
This identification is useful conceptually, but the explicit transition
presentation will be retained because it keeps the Mealy orientation visible.

\section{Input and output labels}
\label{subsec:execution-labels}

Each execution arrow carries two categorical labels. The input label records
which arrow of \(\mathbb A\) is processed. The output label records the
comparison arrow produced by the Mealy output operation.

\begin{definition}[Input labelling functor]
\label{def:input-labelling-functor}
The \textbf{input labelling functor}
\[
    I_\Phi:
    \operatorname{Exec}(\Phi)
    \longrightarrow
    \mathbb A^{\mathrm{op}}
\]
has object part
\[
    p_P:P\to A_0
\]
and arrow part
\[
    (a,x)\longmapsto a^{\mathrm{op}}.
\]
\end{definition}

For \(a:u\to v\), the execution arrow
\[
    (a,x):x\to\delta_P(a,x)
\]
lies over
\[
    a^{\mathrm{op}}:v\to u.
\]
The reversal is forced by the target-to-source processing convention.

\begin{definition}[Output-comparison labelling functor]
\label{def:output-labelling-functor}
The \textbf{output-comparison labelling functor}
\[
    O_\Phi:
    \operatorname{Exec}(\Phi)
    \longrightarrow
    \mathbb B^{\mathrm{op}}
\]
has object part
\[
    q_P:P\to B_0
\]
and arrow part
\[
    (a,x)
    \longmapsto
    \omega_P(a,x)^{\mathrm{op}}.
\]
\end{definition}

The execution category therefore sits in a span
\[
\begin{tikzcd}[column sep=large]
    \mathbb A^{\mathrm{op}}
    &
    \operatorname{Exec}(\Phi)
        \arrow[l, "I_\Phi"']
        \arrow[r, "O_\Phi"]
    &
    \mathbb B^{\mathrm{op}}.
\end{tikzcd}
\]
This is the one-step labelled execution semantics of \(\Phi\).

\begin{proposition}[Functoriality of the execution labels]
\label{prop:functoriality-execution-labels}
The displayed maps define internal functors
\[
    I_\Phi:
    \operatorname{Exec}(\Phi)\to\mathbb A^{\mathrm{op}},
\]
\[
    O_\Phi:
    \operatorname{Exec}(\Phi)\to\mathbb B^{\mathrm{op}}.
\]
\end{proposition}

\begin{proof}
The input labelling preserves identity arrows because the input label of the
identity execution at \(x\) is
\[
    e_A(p_Px)^{\mathrm{op}}.
\]
For composable execution arrows labelled by \(b^{\mathrm{op}}\) and
\(a^{\mathrm{op}}\), one has
\[
    (b\circ a)^{\mathrm{op}}
    =
    a^{\mathrm{op}}\circ b^{\mathrm{op}},
\]
which is exactly the composite in \(\mathbb A^{\mathrm{op}}\) in execution
order.

For output labels, the Mealy output unit equation gives
\[
    \omega_P(e_A(p_Px),x)
    =
    e_B(q_Px),
\]
so identities are preserved. For a composable pair, the output multiplication
law is
\[
    \omega_P(b\circ a,x)
    =
    \omega_P(b,x)
    \circ
    \omega_P\bigl(a,\delta_P(b,x)\bigr).
\]
Passing to \(\mathbb B^{\mathrm{op}}\) gives
\[
    \omega_P(b\circ a,x)^{\mathrm{op}}
    =
    \omega_P\bigl(a,\delta_P(b,x)\bigr)^{\mathrm{op}}
    \circ
    \omega_P(b,x)^{\mathrm{op}},
\]
which is preservation of execution composition.
\end{proof}

\begin{theorem}[Labelled execution category]
\label{thm:labelled-execution-category}
Every Mealy morphism
\[
    \Phi:\mathbb A\rightsquigarrow\mathbb B
\]
determines a span of internal categories
\[
\begin{tikzcd}[column sep=large]
    \mathbb A^{\mathrm{op}}
    &
    \operatorname{Exec}(\Phi)
        \arrow[l, "I_\Phi"']
        \arrow[r, "O_\Phi"]
    &
    \mathbb B^{\mathrm{op}}.
\end{tikzcd}
\]
Its object map records the input and output anchors of a state, and its arrow
map records the input and output-comparison labels of a one-step execution.
\end{theorem}

\section{Finite execution objects}
\label{subsec:finite-execution-objects}

The execution category determines finite runs before any abstract path site is
introduced. For \(n\geq1\), define
\[
    \mathscr T_\Phi(n)
    :=
    \underbrace{
    \operatorname{Exec}(\Phi)_1
    \times_P
    \cdots
    \times_P
    \operatorname{Exec}(\Phi)_1
    }_{n\text{ factors}},
\]
where the pullbacks identify the target of one execution step with the source
of the next. Put
\[
    \mathscr T_\Phi(0):=P,
    \qquad
    \mathscr T_\Phi(-1):=1.
\]

A generalized element of \(\mathscr T_\Phi(n)\) is written
\[
\begin{tikzcd}[column sep=large]
    x_0
        \arrow[r, "{(a_1,x_0)}"]
    &
    x_1
        \arrow[r, "{(a_2,x_1)}"]
    &
    \cdots
        \arrow[r]
    &
    x_n,
\end{tikzcd}
\]
where
\[
    x_i=\delta_P(a_i,x_{i-1})
\]
and
\[
    a_i:p_P(x_i)\to p_P(x_{i-1})
\]
in \(\mathbb A\). The corresponding input path is
\[
\begin{tikzcd}[column sep=large]
    p_P(x_0)
        \arrow[r, "a_1^{\mathrm{op}}"]
    &
    p_P(x_1)
        \arrow[r, "a_2^{\mathrm{op}}"]
    &
    \cdots
        \arrow[r]
    &
    p_P(x_n)
\end{tikzcd}
\]
in \(\mathbb A^{\mathrm{op}}\), while the output-comparison path is
\[
\begin{tikzcd}[column sep=large]
    q_P(x_0)
        \arrow[r, "{\omega_P(a_1,x_0)^{\mathrm{op}}}"]
    &
    q_P(x_1)
        \arrow[r, "{\omega_P(a_2,x_1)^{\mathrm{op}}}"]
    &
    \cdots
        \arrow[r]
    &
    q_P(x_n)
\end{tikzcd}
\]
in \(\mathbb B^{\mathrm{op}}\).

For \(0\leq i\leq n\), write
\[
    \nu_i^{\Phi,n}:
    \mathscr T_\Phi(n)\to P
\]
for the \(i\)-th state map. In particular,
\[
    s_n^\Phi:=\nu_0^{\Phi,n},
    \qquad
    t_n^\Phi:=\nu_n^{\Phi,n}.
\]

For \(0\leq i\leq j\leq n\), projection to the factors between boundaries
\(i\) and \(j\) gives a restriction map
\[
    \operatorname{res}_{i,j}:
    \mathscr T_\Phi(n)
    \longrightarrow
    \mathscr T_\Phi(j-i).
\]
When \(i=j\), this is the vertex map \(\nu_i^{\Phi,n}\). Restriction to the
empty interval is the unique map to \(1\). Nested restrictions compose as
expected because both sides select the same consecutive execution factors.

Adjacent finite executions glue along their common boundary state. For
\(m,n\geq0\), the iterated-pullback definitions give a canonical isomorphism
\[
\begin{tikzcd}[column sep=large]
    \mathscr T_\Phi(m+n)
        \arrow[r, "\cong"]
    &
    \mathscr T_\Phi(m)
        \times_{P}
    \mathscr T_\Phi(n),
\end{tikzcd}
\]
where the pullback uses the terminal boundary of the first execution and the
initial boundary of the second. This is the first appearance of temporal
gluing: compatible adjacent runs determine a unique longer run.

\subsection{Deterministic lifting of finite input paths}
\label{subsec:deterministic-path-lifting}

For \(n\geq1\), let
\[
    \mathbb A^{\mathrm{op}}[n]
    :=
    \underbrace{
    A_1\times_{A_0}\cdots\times_{A_0}A_1
    }_{n\text{ factors}}
\]
denote the object of \(n\)-step paths in \(\mathbb A^{\mathrm{op}}\). Put
\[
    \mathbb A^{\mathrm{op}}[0]:=A_0,
    \qquad
    \mathbb A^{\mathrm{op}}[-1]:=1.
\]
Write
\[
    \lambda_0^{A,n}:
    \mathbb A^{\mathrm{op}}[n]\to A_0
\]
for the initial-boundary map.

The input labelling functor induces
\[
    I_{\Phi,n}:
    \mathscr T_\Phi(n)
    \longrightarrow
    \mathbb A^{\mathrm{op}}[n].
\]
Together with the initial-state map, this gives
\[
\begin{tikzcd}[column sep=large]
    \mathscr T_\Phi(n)
        \arrow[r, "{(s_n^\Phi,I_{\Phi,n})}"]
    &
    P\times_{A_0}\mathbb A^{\mathrm{op}}[n].
\end{tikzcd}
\]

\begin{proposition}[Deterministic path lifting]
\label{prop:deterministic-path-lifting}
For every \(n\geq0\), the canonical map
\[
    \left(
        s_n^\Phi,
        I_{\Phi,n}
    \right):
    \mathscr T_\Phi(n)
    \longrightarrow
    P\times_{A_0}\mathbb A^{\mathrm{op}}[n]
\]
is an isomorphism. The pullback uses
\[
    p_P:P\to A_0
\]
and
\[
    \lambda_0^{A,n}:
    \mathbb A^{\mathrm{op}}[n]\to A_0.
\]
\end{proposition}

\begin{proof}
A finite execution determines its initial state and its input path, so the
displayed map is canonical.

Conversely, suppose internally that one is given an initial state \(x_0\)
and an input path
\[
\begin{tikzcd}[column sep=large]
    p_P(x_0)
        \arrow[r, "a_1^{\mathrm{op}}"]
    &
    u_1
        \arrow[r, "a_2^{\mathrm{op}}"]
    &
    \cdots
        \arrow[r]
    &
    u_n
\end{tikzcd}
\]
in \(\mathbb A^{\mathrm{op}}\). Equivalently, one has composable arrows
\[
    a_i:u_i\to u_{i-1}
\]
in \(\mathbb A\), with \(u_0=p_P(x_0)\). Define recursively
\[
    x_i:=\delta_P(a_i,x_{i-1}).
\]
The typing equations for \(\delta_P\) give
\[
    p_P(x_i)=u_i,
\]
so every successive transition is defined. This produces a finite execution
over the prescribed input path.

Internally, the recursive construction is a morphism in \(\mathcal E\): it is
obtained by repeated use of \(\delta_P\) and the pullback projections of the
input-path object. The construction is unique because the transition is a
map, rather than a relation. It is inverse to the canonical map above.
\end{proof}

\begin{corollary}[Finite executions from states and input paths]
\label{cor:executions-input-paths-initial-states}
At every finite path length, an execution is uniquely determined by
\begin{enumerate}
    \item its initial state;
    \item its input path.
\end{enumerate}
The complete state path and output-comparison path are then forced by the
Mealy transition and output operations.
\end{corollary}

The output path is not additional nondeterministic data. Under the isomorphism
of Proposition~\ref{prop:deterministic-path-lifting}, it is the morphism
obtained by evaluating \(\omega_P\) successively at the recursively generated
states.

\section{Restriction, identity insertion, contraction, and segmentation}
\label{subsec:execution-identity-contraction}

Finite execution objects have restriction and gluing because they are path
objects of the underlying execution graph. The category structure of
\(\operatorname{Exec}(\Phi)\) supplies additional operations.

First, an identity execution may be inserted at any temporal boundary. At a
state \(x_i\), the inserted step is
\[
\begin{tikzcd}[column sep=huge]
    x_i
        \arrow[r, "{(e_A(p_Px_i),x_i)}"]
    &
    x_i.
\end{tikzcd}
\]
Second, adjacent steps may be contracted. For a two-step execution
\[
\begin{tikzcd}[column sep=large]
    x
        \arrow[r, "{(b,x)}"]
    &
    \delta_P(b,x)
        \arrow[r, "{(a,\delta_P(b,x))}"]
    &
    \delta_P\bigl(a,\delta_P(b,x)\bigr),
\end{tikzcd}
\]
contraction produces
\[
\begin{tikzcd}[column sep=huge]
    x
        \arrow[r, "{(b\circ a,x)}"]
    &
    \delta_P(b\circ a,x).
\end{tikzcd}
\]
The transition multiplication law identifies the terminal states, and the
output multiplication law identifies the output label of the contracted step
with the composite of the two original output labels.

More generally, any consecutive block of an \(n\)-step execution may be
contracted to a single step. Total contraction gives a map
\[
    \kappa_n^\Phi:
    \mathscr T_\Phi(n)
    \longrightarrow
    \mathscr T_\Phi(1)
\]
with
\[
    s_1^\Phi\kappa_n^\Phi=s_n^\Phi,
    \qquad
    t_1^\Phi\kappa_n^\Phi=t_n^\Phi.
\]
For \(n=0\), total contraction means insertion of the identity execution.
Associativity says that repeated contractions are independent of
parenthesization, and the unit laws say that inserted identity steps do not
change the contracted execution.

\begin{remark}[Path length records a chosen segmentation]
\label{rem:temporal-extent-segmentation}
The integer \(n\) records a chosen decomposition of an execution into
\(n\) categorical arrows. It does not, in general, measure an irreducible
duration or a number of primitive computational steps.

Indeed, every composable string
\[
    x_0\to x_1\to\cdots\to x_n
\]
contracts to a one-step execution with the same initial and terminal states.
The finite-path objects retain the chosen segmentation, while the contraction
maps record the corresponding categorical composite.

Consequently, endpoint reachability for the full categorical execution
system is already detected by one-step executions: a one-step input label may
itself be an arbitrary composite arrow of \(\mathbb A\).  The later
dynamic-logic chapter will formulate the corresponding universal
finite-trajectory conditions.  To make path length count primitive steps, one
must choose a generating graph or another class of atomic transitions.  The
letterwise Boolean execution graph later in the chapter is the basic example.
\end{remark}

\section{Functoriality under strict semantic homomorphisms}
\label{subsec:temporal-functoriality-strict}

Let
\[
    f:\Phi\to\Psi
\]
be a strict semantic homomorphism between
\[
    \Phi=(P,\delta_P,\omega_P),
    \qquad
    \Psi=(Q,\delta_Q,\omega_Q),
\]
over the same input and output interfaces. Thus the carrier map
\[
    f:P\to Q
\]
preserves both anchors, transitions, and output comparisons.

Define
\[
    \operatorname{Exec}(f):
    \operatorname{Exec}(\Phi)
    \longrightarrow
    \operatorname{Exec}(\Psi)
\]
by
\[
    x\longmapsto f(x)
\]
on objects and
\[
    (a,x)\longmapsto(a,f(x))
\]
on arrows.

\begin{proposition}[Functoriality of execution categories]
\label{prop:functoriality-execution-categories}
The displayed maps define an internal functor
\[
    \operatorname{Exec}(f):
    \operatorname{Exec}(\Phi)
    \to
    \operatorname{Exec}(\Psi).
\]
Moreover,
\[
    I_\Psi\operatorname{Exec}(f)=I_\Phi,
    \qquad
    O_\Psi\operatorname{Exec}(f)=O_\Phi.
\]
\end{proposition}

\begin{proof}
Strict transition preservation gives
\[
    f\bigl(\delta_P(a,x)\bigr)
    =
    \delta_Q(a,f(x)),
\]
so the arrow map preserves targets. It plainly preserves sources. Identity
and composition preservation follow from the formula and the fact that the
input arrow is unchanged.

Input labels are therefore unchanged. Strict output preservation gives
\[
    \omega_Q(a,f(x))
    =
    \omega_P(a,x),
\]
so the output-comparison labels are unchanged as well.
\end{proof}

The induced maps on path objects are
\[
    \mathscr T_f(n):
    \mathscr T_\Phi(n)
    \longrightarrow
    \mathscr T_\Psi(n),
\]
obtained by applying \(f\) at every state boundary and leaving the input and
output labels fixed. These maps commute with contiguous restriction, gluing,
identity insertion, and contraction.

\begin{remark}[Strict homomorphisms and displayed simulations]
\label{rem:temporal-strict-versus-simulation}
The comparison notion in this section is deliberately strict.  A strict
semantic homomorphism has a covariant state map
\[
    f:P\to Q
\]
that preserves transitions and output comparisons on the nose, and therefore
induces a covariant execution functor
\[
    \operatorname{Exec}(\Phi)\to\operatorname{Exec}(\Psi).
\]

A displayed Mealy simulation has different variance.  Its displayed state
component is directed from the target carrier to the source carrier and its
output component is a comparison arrow rather than a strict equality of
outputs.  Such a simulation does not directly define a labelled functor
between the unevaluated execution categories considered here.  In the next
chapter, the output-comparison arrow is evaluated in a downstream
\(\mathbb B\)-action.  That evaluation produces the contravariant relative
execution functor associated to the simulation.  Thus the covariant
functoriality of the present section and the contravariant functoriality of
evaluated simulations are complementary results for two genuinely different
comparison notions.
\end{remark}

\section{Composition of execution categories}
\label{sec:composition-temporal-execution}

The temporal semantics of a composite Mealy morphism is obtained by matching
the output-comparison execution of the first component with the input
execution of the second component.  This section makes that matching precise
as a pullback over the intermediate labelled category.

\subsection{Composite Mealy morphisms}
\label{subsec:composite-mealy-morphisms}

Let
\[
    \Phi=(P,\delta_P,\omega_P):
    \mathbb A\rightsquigarrow\mathbb B
\]
and
\[
    \Psi=(Q,\delta_Q,\omega_Q):
    \mathbb B\rightsquigarrow\mathbb C
\]
be composable Mealy morphisms. Write
\[
    p_Q:Q\to B_0
\]
for the input anchor of \(Q\). The carrier of the composite is
\[
    P\odot Q
    :=
    P\times_{q_P,B_0,p_Q}Q.
\]

For a state pair \((x,y)\) and an input arrow \(a\) applicable to \(x\), put
\[
    b:=\omega_P(a,x).
\]
Since
\[
    t_B(b)=q_P(x)=p_Q(y),
\]
the arrow \(b\) is applicable to \(y\). The transition and output comparison
of the composite are
\[
    \delta_{\Psi\circ\Phi}(a,(x,y))
    =
    \left(
        \delta_P(a,x),
        \delta_Q(b,y)
    \right),
\]
\[
    \omega_{\Psi\circ\Phi}(a,(x,y))
    =
    \omega_Q(b,y).
\]
Thus one step of the composite consists of one step of \(\Phi\), whose output
comparison is immediately processed as the input arrow for one step of
\(\Psi\).

\subsection{Execution categories compose by pullback}
\label{subsec:execution-category-pullback}

The shared interface is represented by the functors
\[
    O_\Phi:
    \operatorname{Exec}(\Phi)\to\mathbb B^{\mathrm{op}},
\]
\[
    I_\Psi:
    \operatorname{Exec}(\Psi)\to\mathbb B^{\mathrm{op}}.
\]
Their pullback pairs executions whose intermediate labels agree.

\begin{theorem}[Execution-category composition]
\label{thm:execution-category-composition}
There is a canonical isomorphism of internal categories
\[
\begin{tikzcd}[column sep=large]
    \operatorname{Exec}(\Psi\circ\Phi)
        \arrow[r, "\cong"]
    &
    \operatorname{Exec}(\Phi)
        \times_{\mathbb B^{\mathrm{op}}}
    \operatorname{Exec}(\Psi).
\end{tikzcd}
\]
Under this isomorphism, the input labelling is induced by \(I_\Phi\), and the
output-comparison labelling is induced by \(O_\Psi\).
\end{theorem}

\begin{proof}
On objects, the pullback is
\[
    P\times_{B_0}Q,
\]
which is the carrier of \(\Psi\circ\Phi\).

Define the arrow component of the comparison by
\[
    (a,x,y)
    \longmapsto
    \left(
        (a,x),
        (\omega_P(a,x),y)
    \right).
\]
The second execution arrow is defined because
\[
    t_B(\omega_P(a,x))
    =
    q_P(x)
    =
    p_Q(y).
\]
Moreover, the two arrows have the same label in
\(\mathbb B^{\mathrm{op}}\):
\[
\begin{tikzcd}[column sep=large]
    (a,x)
        \arrow[d, mapsto, "O_\Phi"']
    &
    (\omega_P(a,x),y)
        \arrow[d, mapsto, "I_\Psi"]
    \\
    \omega_P(a,x)^{\mathrm{op}}
        \arrow[r, equal]
    &
    \omega_P(a,x)^{\mathrm{op}}.
\end{tikzcd}
\]

Conversely, an arrow of the pullback consists of an execution arrow
\((a,x)\) of \(\Phi\), an execution arrow \((b,y)\) of \(\Psi\), and an
equality
\[
    b^{\mathrm{op}}
    =
    \omega_P(a,x)^{\mathrm{op}}
\]
in \(\mathbb B^{\mathrm{op}}\). Equivalently,
\[
    b=\omega_P(a,x).
\]
Hence the pullback arrow is uniquely determined by \((a,x,y)\), exactly the
data of a one-step execution of the composite.

Its target is
\[
    \left(
        \delta_P(a,x),
        \delta_Q(\omega_P(a,x),y)
    \right),
\]
which is the composite transition. Identities agree because
\[
    \omega_P(e_A(p_Px),x)=e_B(q_Px).
\]

For composition, let
\[
    a:u\to v,
    \qquad
    b:v\to w,
\]
and let \(x\) lie over \(w\). Put
\[
    x_1:=\delta_P(b,x),
\]
\[
    \beta_1:=\omega_P(b,x),
    \qquad
    \beta_2:=\omega_P(a,x_1).
\]
The output multiplication law gives
\[
    \omega_P(b\circ a,x)
    =
    \beta_1\circ\beta_2.
\]
The transition multiplication law for \(\Psi\) gives
\[
    \delta_Q(\beta_1\circ\beta_2,y)
    =
    \delta_Q\bigl(\beta_2,\delta_Q(\beta_1,y)\bigr).
\]
Therefore composition in the pullback agrees with composition in
\(\operatorname{Exec}(\Psi\circ\Phi)\). The objectwise and arrowwise
comparisons are mutually inverse internal functors.
\end{proof}

\subsection{Naturality, units, and coherence}
\label{subsec:naturality-units-coherence}

Let
\[
    f:\Phi\to\Phi',
    \qquad
    g:\Psi\to\Psi'
\]
be strict semantic homomorphisms. Their horizontal composite has carrier map
\[
    f\odot g:
    P\times_{B_0}Q
    \longrightarrow
    P'\times_{B_0}Q',
    \qquad
    (x,y)\longmapsto(f(x),g(y)).
\]

\begin{proposition}[Naturality of execution-category composition]
\label{prop:naturality-execution-composition}
The isomorphism of
Theorem~\ref{thm:execution-category-composition}
is natural in strict semantic homomorphisms of both Mealy factors.
\end{proposition}

\begin{proof}
Let
\[
    \kappa_{\Phi,\Psi}:
    \operatorname{Exec}(\Psi\circ\Phi)
    \xrightarrow{\cong}
    \operatorname{Exec}(\Phi)
    \times_{\mathbb B^{\mathrm{op}}}
    \operatorname{Exec}(\Psi)
\]
denote the comparison. On a one-step composite execution,
\[
    \kappa_{\Phi,\Psi}(a,x,y)
    =
    \left(
        (a,x),
        (\omega_\Phi(a,x),y)
    \right).
\]
Applying
\(\operatorname{Exec}(f)\times\operatorname{Exec}(g)\) gives
\[
    \left(
        (a,f(x)),
        (\omega_\Phi(a,x),g(y))
    \right).
\]
Strict output preservation gives
\[
    \omega_{\Phi'}(a,f(x))
    =
    \omega_\Phi(a,x),
\]
so this is exactly
\[
    \kappa_{\Phi',\Psi'}(a,f(x),g(y)).
\]
The comparison square therefore commutes on objects and arrows.
\end{proof}

Let
\[
    1_{\mathbb A}:
    \mathbb A\rightsquigarrow\mathbb A
\]
be the identity Mealy morphism. Its carrier is \(A_0\), and its structure is
\[
    \delta(a,t_Aa)=s_Aa,
    \qquad
    \omega(a,t_Aa)=a.
\]

\begin{proposition}[Execution unit]
\label{prop:temporal-unit}
There is a canonical isomorphism
\[
    \operatorname{Exec}(1_{\mathbb A})
    \cong
    \mathbb A^{\mathrm{op}},
\]
compatible with both input and output-comparison labelling.
\end{proposition}

\begin{proof}
The object object is \(A_0\), and
\[
    A_1\times_{t_A,A_0}A_0
    \cong
    A_1.
\]
Under this identification, the source and target maps of the execution
category are \(t_A\) and \(s_A\), respectively, so the execution category is
\(\mathbb A^{\mathrm{op}}\). Both labels send the execution represented by
\(a\) to \(a^{\mathrm{op}}\), and the identity and composition structures
agree.
\end{proof}

\begin{proposition}[Associative coherence]
\label{prop:temporal-associative-coherence}
The execution-category composition isomorphisms are compatible with the
associativity and unit isomorphisms for Mealy composition.
\end{proposition}

\begin{proof}
For three composable Mealy morphisms, both parenthesized pullbacks of execution
categories are obtained from the same iterated pullback over the two shared
interfaces. Every route around the associativity pentagon has the same
projections to the three execution-category factors, and every route around a
unit triangle has the same projection to the remaining factor. Equality
follows from the uniqueness clause of the pullback universal property.
\end{proof}

\begin{remark}[Execution spans form a pseudofunctor]
\label{rem:execution-pseudofunctor}
The assignments
\[
    \mathbb A\longmapsto\mathbb A^{\mathrm{op}},
\]
\[
    \Phi\longmapsto
    \left(
        \mathbb A^{\mathrm{op}}
        \xleftarrow{I_\Phi}
        \operatorname{Exec}(\Phi)
        \xrightarrow{O_\Phi}
        \mathbb B^{\mathrm{op}}
    \right),
\]
and
\[
    f\longmapsto\operatorname{Exec}(f)
\]
define a pseudofunctor, equivalently a homomorphism of bicategories, from the
bicategory of Mealy morphisms and strict semantic homomorphisms to the
bicategory of spans of internal categories in \(\mathcal E\).
\end{remark}

\begin{remark}[Forward specialization to downstream evaluation]
\label{rem:temporal-downstream-specialization}
Let \(Y\to B_0\) be a right \(\mathbb B\)-action.  The next chapter packages
\(Y\) as a terminal-output Mealy morphism
\[
    \widehat Y:\mathbb B\rightsquigarrow\mathbf 1
\]
whose execution category is the action category
\[
    \operatorname{Exec}(\widehat Y)
    \cong
    \int_{\mathbb B}Y.
\]
The pullback theorem of this section then specializes to
\[
    \operatorname{Exec}(\widehat Y\circ\Phi)
    \cong
    \operatorname{Exec}(\Phi)
    \times_{\mathbb B^{\mathrm{op}}}
    \int_{\mathbb B}Y.
\]
The left-hand side will be identified there with the relative program
category of \(\Phi\) evaluated in \(Y\).  No new execution-category laws are
required: they are inherited from Mealy composition and the pullback
description proved above.
\end{remark}

\section{The finite-path site and trajectory sheaves}
\label{sec:internal-graphs-trajectory-sheaves}

The construction of \(\mathscr T_\Phi(n)\) used only the underlying graph of
\(\operatorname{Exec}(\Phi)\): its state object, its one-step execution
object, and the source and target maps. Contiguous restriction and gluing were
already present before identity insertion and contraction were considered.
This section isolates that graph-level structure.

\subsection{The augmented discrete interval category}
\label{subsec:augmented-discrete-interval-category}

Let
\[
    [-1]:=\varnothing
\]
denote the empty finite ordinal. For \(n\geq0\), write
\[
    [n]:=\{0<1<\cdots<n\}.
\]
The object \([n]\) represents a finite interval with \(n\) one-step segments
and \(n+1\) boundary instants.

\begin{definition}[Discrete interval category]
\label{def:discrete-interval-category}
The \textbf{augmented discrete interval category}
\[
    \mathbf{Int}_{\mathrm{disc}}
\]
has the finite ordinals
\[
    [-1],[0],[1],[2],\ldots
\]
as objects. A morphism
\[
    \iota:[k]\longrightarrow[n]
\]
is an injective order-preserving map whose image is a contiguous subinterval
of \([n]\). Equivalently, for \(k\geq0\), it is determined by a starting
position \(i\) with
\[
    0\leq i\leq n-k
\]
and has the form
\[
    \iota(r)=i+r.
\]
There is a unique morphism
\[
    [-1]\to[n]
\]
for every \(n\geq-1\).
\end{definition}

In starting-position notation,
\[
    (i,k):[k]\to[n]
\]
and
\[
    (j,\ell):[\ell]\to[k]
\]
compose to
\[
    (i+j,\ell):[\ell]\to[n].
\]
The two maps
\[
    \epsilon_0,\epsilon_1:[0]\rightrightarrows[1]
\]
select the initial and terminal boundaries. For \(n\geq1\) and
\(0\leq i<n\), write
\[
    e_i:[1]\to[n]
\]
for the elementary edge inclusion with image \(\{i,i+1\}\).

Intersections of contiguous subintervals are contiguous, possibly empty.
Consequently, pullbacks of interval inclusions exist in
\(\mathbf{Int}_{\mathrm{disc}}\). For example, adjacent elementary edges
intersect in a point,
\[
\begin{tikzcd}
    {[0]}
        \arrow[r]
        \arrow[d]
    &
    {[1]}
        \arrow[d, "e_{i+1}"]
    \\
    {[1]}
        \arrow[r, "e_i"']
    &
    {[n]},
\end{tikzcd}
\]
whereas nonadjacent elementary edges intersect in \([-1]\).

\subsection{The elementary-edge topology}
\label{subsec:path-topology}

\begin{definition}[Elementary path covers]
\label{def:elementary-path-covers}
For \(n\geq1\), the \textbf{elementary path cover} of \([n]\) is the family
\[
    \{e_i:[1]\to[n]\}_{0\leq i<n}.
\]
The identity family covers \([0]\), and the empty family covers \([-1]\).
\end{definition}

\begin{definition}[Finite-path topology]
\label{def:discrete-path-topology}
The \textbf{finite-path topology}
\[
    J_{\mathrm{path}}
\]
is the Grothendieck topology on \(\mathbf{Int}_{\mathrm{disc}}\) generated by
the elementary path covers, the identity cover of \([0]\), and the empty
cover of \([-1]\).
\end{definition}

We verify that the displayed covering families form a basis. Let
\[
    \iota:[k]\to[n]
\]
be a contiguous interval inclusion. The pullback of an elementary edge
inclusion along \(\iota\) is an elementary edge, a boundary point, or the
empty interval.

Suppose first that \(k\geq1\) and that the image of \(\iota\) is
\[
    \{i,i+1,\ldots,i+k\}.
\]
For every \(0\leq r<k\), the pullback of
\[
    e_{i+r}:[1]\to[n]
\]
along \(\iota\) is the elementary edge
\[
    e_r:[1]\to[k].
\]
Hence the pullback family contains the full elementary-edge cover of
\([k]\), possibly together with boundary-point and empty intersections. If
\(k=0\), the pullback family contains the identity of \([0]\). If
\(k=-1\), every pullback intersection is \([-1]\), and the resulting family
is refined by the specified empty covering family of \([-1]\).

Transitivity is equally direct. Every member of an elementary-edge cover is a
copy of \([1]\), whose elementary cover is its identity family. Refining
each edge by that family reproduces the original cover. The identity and
empty covers are similarly stable. Thus the displayed families form a basis
for \(J_{\mathrm{path}}\).

\subsection{Trajectory presheaves and their gluing condition}
\label{subsec:concrete-path-sheaf-condition}

\begin{definition}[Trajectory presheaf]
\label{def:internal-trajectory-presheaf}
An \(\mathcal E\)-valued \textbf{trajectory presheaf} is a functor
\[
    X:
    \mathbf{Int}_{\mathrm{disc}}^{\mathrm{op}}
    \to
    \mathcal E.
\]
For a contiguous inclusion
\[
    \iota:[k]\to[n],
\]
the induced map
\[
    X(\iota):X(n)\to X(k)
\]
is called contiguous path restriction.
\end{definition}

\begin{definition}[Trajectory sheaf]
\label{def:discrete-trajectory-sheaf}
A trajectory presheaf is a \textbf{trajectory sheaf} if it is a sheaf for
\(J_{\mathrm{path}}\).
\end{definition}

Write
\[
    \mathbf{TrajSh}_{\mathcal E}
    :=
    \mathbf{Sh}_{\mathcal E}
    \left(
        \mathbf{Int}_{\mathrm{disc}},
        J_{\mathrm{path}}
    \right)
\]
for the category of trajectory sheaves and their morphisms.

For the elementary cover of \([n]\), a matching family consists of one
section of \(X(1)\) for each elementary edge. Adjacent sections must agree on
their common boundary point, and nonadjacent sections must agree on their
empty intersection. The latter conditions become automatic once
\(X(-1)\cong1\).

\begin{proposition}[Concrete trajectory-sheaf condition]
\label{prop:concrete-trajectory-sheaf-condition}
A trajectory presheaf \(X\) is a trajectory sheaf if and only if:
\begin{enumerate}
    \item the canonical map
    \[
        X(-1)\longrightarrow1
    \]
    is an isomorphism;
    \item for every \(n\geq1\), the canonical edge-restriction map
    \[
        X(n)
        \longrightarrow
        \underbrace{
        X(1)\times_{X(0)}
        X(1)\times_{X(0)}
        \cdots\times_{X(0)}X(1)
        }_{n\text{ factors}}
    \]
    is an isomorphism.
\end{enumerate}
The pullbacks use terminal restriction on one edge and initial restriction on
the next.
\end{proposition}

\begin{proof}
The empty covering family of \([-1]\) has terminal matching object. Hence the
sheaf condition at \([-1]\) is precisely
\[
    X(-1)\cong1.
\]

For \(n\geq1\), a matching family for the elementary-edge cover consists
internally of sections
\[
    x_i\in X(1),
    \qquad
    0\leq i<n,
\]
satisfying
\[
    X(\epsilon_1)(x_i)
    =
    X(\epsilon_0)(x_{i+1})
\]
for every adjacent pair. These are exactly the equations defining the
iterated pullback
\[
    X(1)\times_{X(0)}\cdots\times_{X(0)}X(1).
\]
Intersections of nonadjacent edges are empty, and compatibility over them is
automatic because \(X(-1)\cong1\). The sheaf comparison for the basis cover
is therefore the displayed edge-restriction map. Since the specified covers
form a basis, the basis sheaf condition is equivalent to the full sheaf
condition.
\end{proof}

\begin{corollary}[Adjacent path gluing]
\label{cor:adjacent-temporal-gluing}
For every trajectory sheaf \(X\) and every \(m,n\in\mathbb N\), there is a
canonical isomorphism
\[
\begin{tikzcd}[column sep=large]
    X(m+n)
        \arrow[r, "\cong"]
    &
    X(m)\times_{X(0)}X(n),
\end{tikzcd}
\]
where the first map to \(X(0)\) is restriction to the terminal boundary of
\([m]\), and the second is restriction to the initial boundary of \([n]\).
\end{corollary}

\begin{proof}
Both sides are canonically isomorphic to the iterated pullback of \(m+n\)
copies of \(X(1)\) over \(X(0)\). The first \(m\) factors form the initial
path, the final \(n\) factors form the terminal path, and the middle pullback
equation identifies their shared boundary. The cases \(m=0\) or \(n=0\)
reduce to pullback along an identity map.
\end{proof}

\subsection{Path objects of an internal graph}
\label{subsec:path-objects-internal-graph}

Let
\[
    G=
    \left(
    G_1
    \mathrel{\substack{\xrightarrow{t_G}\\[-0.4em]\xrightarrow[s_G]{}}}
    G_0
    \right)
\]
be an internal graph in \(\mathcal E\).

\begin{definition}[Graph path objects]
\label{def:graph-path-objects}
Define
\[
    G^{[-1]}:=1,
    \qquad
    G^{[0]}:=G_0,
    \qquad
    G^{[1]}:=G_1.
\]
For \(n\geq2\), define
\[
    G^{[n]}
    :=
    \underbrace{
    G_1\times_{G_0}G_1
    \times_{G_0}\cdots
    \times_{G_0}G_1
    }_{n\text{ factors}},
\]
where the target of one edge is identified with the source of the next.
\end{definition}

A generalized element of \(G^{[n]}\) is a graph path
\[
\begin{tikzcd}[column sep=large]
    x_0
        \arrow[r, "g_1"]
    &
    x_1
        \arrow[r, "g_2"]
    &
    \cdots
        \arrow[r, "g_n"]
    &
    x_n.
\end{tikzcd}
\]
For \(0\leq i\leq n\), write
\[
    \nu_i^{(n)}:G^{[n]}\to G_0
\]
for the \(i\)-th vertex map.

A contiguous inclusion
\[
    \iota:[k]\to[n]
\]
with image \(\{i,i+1,\ldots,i+k\}\) selects the corresponding subpath. If
\(k=0\), it selects the vertex \(x_i\); restriction to \([-1]\) is the
unique map to \(1\).

\begin{definition}[Graph trajectory sheaf]
\label{def:graph-trajectory-presheaf}
The \textbf{graph trajectory presheaf} of \(G\) is
\[
    \operatorname{Path}(G):
    \mathbf{Int}_{\mathrm{disc}}^{\mathrm{op}}
    \to
    \mathcal E
\]
defined by
\[
    \operatorname{Path}(G)(n):=G^{[n]}.
\]
Contiguous restriction is subpath selection.
\end{definition}

\begin{lemma}[Functoriality of graph paths]
\label{lem:functoriality-graph-trajectories}
The displayed path objects and restriction maps define a functor
\[
    \operatorname{Path}(G):
    \mathbf{Int}_{\mathrm{disc}}^{\mathrm{op}}
    \to
    \mathcal E.
\]
\end{lemma}

\begin{proof}
Restriction along an identity inclusion is the identity map. Restricting to a
contiguous interval and then to a contiguous interval inside it selects the
same edge and vertex factors as restriction along the composite inclusion.
The unique restriction to \([-1]\) is compatible with every composite because
its codomain is terminal.
\end{proof}

\begin{theorem}[Temporal gluing of graph paths]
\label{thm:temporal-gluing-graph-trajectories}
For every internal graph \(G\), the presheaf
\[
    \operatorname{Path}(G)
\]
is a trajectory sheaf.
\end{theorem}

\begin{proof}
The value at the empty interval is
\[
    \operatorname{Path}(G)(-1)=1.
\]
For \(n\geq1\), the elementary-edge comparison is
\[
    G^{[n]}
    \longrightarrow
    \underbrace{
    G_1\times_{G_0}\cdots\times_{G_0}G_1
    }_{n\text{ factors}}.
\]
By definition, the domain is precisely the displayed iterated pullback and
the comparison is its canonical identity comparison. Proposition
\ref{prop:concrete-trajectory-sheaf-condition} therefore applies.
\end{proof}

\subsection{Reconstruction of the graph from its path sheaf}
\label{subsec:reconstruction-underlying-graph}

Let \(X\) be a trajectory sheaf. Define an internal graph \(G_X\) by
\[
    (G_X)_0:=X(0),
    \qquad
    (G_X)_1:=X(1),
\]
with source and target
\[
    s_X:=X(\epsilon_0):X(1)\to X(0),
\]
\[
    t_X:=X(\epsilon_1):X(1)\to X(0).
\]

\begin{proposition}[Reconstruction from lengths zero and one]
\label{prop:reconstruction-length-zero-one}
For every trajectory sheaf \(X\), there are canonical isomorphisms
\[
    \theta_{X,n}:
    X(n)
    \xrightarrow{\cong}
    (G_X)^{[n]}
\]
for all \(n\geq-1\). These isomorphisms are natural in contiguous interval
inclusions.
\end{proposition}

\begin{proof}
For \(n=-1\), the claim is the sheaf isomorphism \(X(-1)\cong1\). For
\(n=0,1\), it follows from the definition of \(G_X\). For \(n\geq2\), the
concrete sheaf condition gives
\[
    X(n)
    \cong
    \underbrace{
    X(1)\times_{X(0)}\cdots\times_{X(0)}X(1)
    }_{n\text{ factors}},
\]
which is exactly \((G_X)^{[n]}\).

The comparison is induced by elementary-edge restriction. A positive-length
subinterval selects the same consecutive edge factors on both sides. A point
subinterval selects the corresponding vertex. At an interior vertex, this
may be computed either as the target of the preceding edge or as the source
of the following edge, and the two maps agree by the defining pullback
equations. Restriction to \([-1]\) is the unique map to \(1\). Hence the
family \(\theta_{X,n}\) is natural in every contiguous inclusion.
\end{proof}

A morphism of trajectory sheaves
\[
    \alpha:X\to Y
\]
has components
\[
    \alpha_0:X(0)\to Y(0),
    \qquad
    \alpha_1:X(1)\to Y(1).
\]
Naturality at \(\epsilon_0\) and \(\epsilon_1\) gives
\[
    s_Y\alpha_1=\alpha_0s_X,
    \qquad
    t_Y\alpha_1=\alpha_0t_X,
\]
so \((\alpha_0,\alpha_1)\) is a graph morphism. Conversely, a graph morphism
induces maps between every iterated pullback of composable edges, and those
maps commute with contiguous restriction.

\begin{theorem}[Graph--trajectory-sheaf equivalence]
\label{thm:graph-trajectory-sheaf-equivalence}
The graph path construction defines an equivalence
\[
\begin{tikzcd}[column sep=large]
    \mathbf{Graph}_{\mathrm{int}}(\mathcal E)
        \arrow[r, shift left=0.7ex, "{\operatorname{Path}}"]
    &
    \mathbf{TrajSh}_{\mathcal E}
        \arrow[l, shift left=0.7ex, "{X\mapsto G_X}"] .
\end{tikzcd}
\]
\end{theorem}

\begin{proof}
For an internal graph \(G\), reconstruction from
\[
    \operatorname{Path}(G)(0)=G_0,
    \qquad
    \operatorname{Path}(G)(1)=G_1
\]
returns \(G\) strictly. For a trajectory sheaf \(X\), Proposition
\ref{prop:reconstruction-length-zero-one} gives a natural isomorphism
\[
    X\cong\operatorname{Path}(G_X).
\]
The preceding discussion shows that the two constructions are mutually
inverse on morphisms. These comparisons supply the unit and counit of an
equivalence.
\end{proof}

\begin{corollary}[Path data of an internal graph]
\label{cor:temporal-data-of-graph}
A trajectory sheaf contains no independent graph-level data in lengths greater
than one. Every finite path object is reconstructed from
\begin{enumerate}
    \item the boundary-state object \(X(0)\);
    \item the one-step object \(X(1)\);
    \item the two endpoint maps
    \[
        X(1)\rightrightarrows X(0).
    \]
\end{enumerate}
\end{corollary}

Applied to the underlying graph of \(\operatorname{Exec}(\Phi)\), the
construction recovers exactly the finite execution objects already defined:
\[
\begin{tikzcd}[column sep=large]
    \operatorname{Path}
    \bigl(U\operatorname{Exec}(\Phi)\bigr)(n)
        \arrow[r, "\cong"]
    &
    \mathscr T_\Phi(n).
\end{tikzcd}
\]
Thus the path site has been extracted from the automata-theoretic execution
construction rather than imposed on it in advance.

\section{Segal composition and internal categories}
\label{sec:segal-enhancement-categorical-execution}

The graph-level path sheaf remembers a string of one-step executions and its
contiguous restrictions. It does not, by itself, specify which arrows are
identities or how adjacent arrows compose. Those operations arise from the
category structure and are encoded by the full simplex category.

\subsection{The augmented simplex category}
\label{subsec:augmented-simplex-category}

Let
\[
    \Delta_+
\]
denote the augmented simplex category with objects
\[
    [-1],[0],[1],[2],\ldots
\]
and all order-preserving maps as morphisms. There is a canonical inclusion
\[
\begin{tikzcd}[column sep=large]
    \mathbf{Int}_{\mathrm{disc}}
        \arrow[r, hook, "j"]
    &
    \Delta_+
\end{tikzcd}
\]
which is the identity on objects and includes the contiguous interval
injections.

The image of \(j\) is the augmented inert subcategory: its nonempty maps are
the distance-preserving injections
\[
    \iota(r)=\iota(0)+r.
\]
In particular, the two outer coface maps
\[
    \delta^0,\delta^n:[n-1]\to[n]
\]
belong to the image of \(j\). Under a contravariant nerve, they give initial
and terminal truncation of paths.

The additional simplicial maps provide the category operations:
\begin{enumerate}
    \item a codegeneracy duplicates a vertex and therefore inserts an
    identity arrow;
    \item an inner coface omits an interior vertex and therefore contracts the
    two adjacent arrows by composition.
\end{enumerate}
The simplicial identities encode the unit and associativity equations among
these operations.

\subsection{The augmented internal nerve}
\label{subsec:augmented-nerve}

Let
\[
    \mathbb C=(C_0,C_1,s_C,t_C,e_C,m_C)
\]
be an internal category.

\begin{definition}[Augmented internal nerve]
\label{def:augmented-internal-nerve}
The \textbf{augmented nerve} of \(\mathbb C\) is the functor
\[
    N_+\mathbb C:
    \Delta_+^{\mathrm{op}}
    \to
    \mathcal E
\]
defined by
\[
    N_+\mathbb C_{-1}:=1,
    \qquad
    N_+\mathbb C_0:=C_0,
\]
and, for \(n\geq1\),
\[
    N_+\mathbb C_n
    :=
    \underbrace{
    C_1\times_{C_0}\cdots\times_{C_0}C_1
    }_{n\text{ factors}}.
\]
Outer face maps truncate a path, inner face maps compose adjacent arrows, and
degeneracy maps insert identities. Every map to degree \(-1\) is the unique
map to \(1\).
\end{definition}

The restriction of \(N_+\mathbb C\) to the ordinary simplex category is the
usual internal nerve.

\begin{definition}[Augmented strict Segal object]
\label{def:augmented-strict-segal-object}
An \textbf{augmented strict Segal object} in \(\mathcal E\) is a functor
\[
    Y:\Delta_+^{\mathrm{op}}\to\mathcal E
\]
such that
\begin{enumerate}
    \item the canonical map \(Y_{-1}\to1\) is an isomorphism;
    \item for every \(n\geq2\), the canonical Segal map
    \[
        Y_n
        \longrightarrow
        \underbrace{
        Y_1\times_{Y_0}\cdots\times_{Y_0}Y_1
        }_{n\text{ factors}}
    \]
    is an isomorphism.
\end{enumerate}
\end{definition}

\begin{proposition}[Strict Segal isomorphisms]
\label{prop:strict-segal-isomorphisms}
For every internal category \(\mathbb C\), the augmented nerve
\(N_+\mathbb C\) is an augmented strict Segal object.
\end{proposition}

\begin{proof}
The degree \(-1\) object is \(1\). In every degree \(n\geq2\), both sides of
the Segal comparison are the same iterated pullback of \(C_1\) over \(C_0\).
The comparison is induced by the elementary-edge inclusions and is the
canonical identity comparison under these pullback descriptions.
\end{proof}

\subsection{The internal nerve equivalence}
\label{subsec:internal-nerve-equivalence}

Write
\[
    \mathsf{Seg}^{\mathrm{str}}_+(\mathcal E)
\]
for the full subcategory of
\([
    \Delta_+^{\mathrm{op}},\mathcal E
]\)
spanned by augmented strict Segal objects.

\begin{theorem}[Internal nerve equivalence]
\label{thm:internal-nerve-equivalence}
The augmented nerve defines an equivalence
\[
\begin{tikzcd}[column sep=large]
    \mathsf{IntCat}(\mathcal E)
        \arrow[r, shift left=0.7ex, "N_+"]
    &
    \mathsf{Seg}^{\mathrm{str}}_+(\mathcal E)
        \arrow[l, shift left=0.7ex, "R"] .
\end{tikzcd}
\]
\end{theorem}

\begin{proof}
Let
\[
    Y:\Delta_+^{\mathrm{op}}\to\mathcal E
\]
be augmented strict Segal. Define \(R(Y)\) by
\[
    R(Y)_0:=Y_0,
    \qquad
    R(Y)_1:=Y_1.
\]
The source and target maps are induced by the two endpoint inclusions,
\[
    s_Y:=Y(\epsilon_0),
    \qquad
    t_Y:=Y(\epsilon_1).
\]
The unique codegeneracy
\[
    \sigma^0:[1]\to[0]
\]
gives the identity map
\[
    e_Y:=Y(\sigma^0):Y_0\to Y_1.
\]
The inner coface
\[
    \delta^1:[1]\to[2],
    \qquad
    0\mapsto0,
    \quad
    1\mapsto2,
\]
gives
\[
    Y(\delta^1):Y_2\to Y_1.
\]
Transporting across the Segal isomorphism defines composition:
\[
\begin{tikzcd}[column sep=large]
    Y_1\times_{Y_0}Y_1
        \arrow[r, "\cong"']
    &
    Y_2
        \arrow[r, "{Y(\delta^1)}"]
    &
    Y_1.
\end{tikzcd}
\]
The simplicial identities involving \(\sigma^0\), the outer cofaces, and
\(\delta^1\) give the source, target, left-unit, right-unit, and associativity
axioms. Hence \(R(Y)\) is an internal category.

A natural transformation of augmented strict Segal objects is determined in
degrees zero and one by the Segal isomorphisms. Naturality at \(\sigma^0\)
and \(\delta^1\) says precisely that these two maps preserve identities and
composition. Therefore \(R\) acts on morphisms by internal functors.

For an internal category \(\mathbb C\), reconstruction from \(N_+\mathbb C\)
returns \(\mathbb C\), up to the canonical identifications of the chosen
pullbacks. Conversely, the Segal maps give degreewise isomorphisms
\[
    Y_n\xrightarrow{\cong}N_+R(Y)_n
\]
for \(n\geq0\), and the augmentation gives
\[
    Y_{-1}\cong1=N_+R(Y)_{-1}.
\]

It remains to check naturality in all simplicial maps. Let
\[
    \varphi:[m]\to[n]
\]
be order preserving. An \(n\)-path is sent to the \(m\)-path whose
\(j\)-th arrow is the composite of the consecutive block lying between the
vertices \(\varphi(j-1)\) and \(\varphi(j)\), with an identity inserted when
those vertices coincide. Every morphism in \(\Delta\) is generated by
cofaces and codegeneracies. Outer cofaces truncate, inner cofaces compose, and
codegeneracies insert identities. The simplicial identities are exactly the
unit and associativity equations showing that the block formula is independent
of a chosen expression in generators. Hence the degreewise Segal
isomorphisms are natural. Maps involving \([-1]\) are automatic because the
degree \(-1\) object is terminal. Thus
\[
    Y\cong N_+R(Y)
\]
as augmented simplicial objects.
\end{proof}

\subsection{The underlying trajectory sheaf of an internal category}
\label{subsec:underlying-trajectory-sheaf}

\begin{definition}[Underlying trajectory sheaf]
\label{def:underlying-trajectory-sheaf}
The \textbf{underlying trajectory sheaf} of an internal category \(\mathbb C\)
is
\[
    \operatorname{Path}(\mathbb C)
    :=
    j^\ast N_+\mathbb C.
\]
\end{definition}

This name emphasizes what the restriction retains: the objects, arrows,
finite composable strings, and contiguous restrictions. It avoids suggesting
that the construction forgets execution data in any semantic sense; it forgets
only the active simplicial operations of identity insertion and composition.

Let \(U\mathbb C\) be the underlying internal graph.

\begin{proposition}[Category paths are graph paths]
\label{prop:category-paths-underlying-graph}
There is a canonical natural isomorphism
\[
\begin{tikzcd}[column sep=large]
    \operatorname{Path}(\mathbb C)
        \arrow[r, "\cong"]
    &
    \operatorname{Path}(U\mathbb C).
\end{tikzcd}
\]
\end{proposition}

\begin{proof}
Both functors take \([n]\) to the object of \(n\) composable arrows of
\(\mathbb C\). A contiguous inclusion restricts such a string to the same
contiguous substring in both constructions. The comparison is therefore the
canonical identity comparison in every degree.
\end{proof}

The underlying trajectory sheaf remembers
\begin{enumerate}
    \item objects and one-step arrows;
    \item finite composable strings;
    \item contiguous subpath restriction;
    \item initial and terminal truncation;
    \item unique gluing along matching boundaries.
\end{enumerate}
The augmented nerve additionally remembers
\begin{enumerate}
    \item which one-step arrows are identities;
    \item how adjacent arrows contract by composition;
    \item the unit and associativity equations governing contraction.
\end{enumerate}

\subsection{Segal composition structures on trajectory sheaves}
\label{subsec:classification-segal-enhancements}

Let \(X\) be a trajectory sheaf. By
Theorem~\ref{thm:graph-trajectory-sheaf-equivalence}, it determines the graph
\[
    G_X=
    \left(
    X(1)
    \mathrel{\substack{\xrightarrow{t_X}\\[-0.4em]\xrightarrow[s_X]{}}}
    X(0)
    \right).
\]
For \(n\geq1\), write
\[
    \pi_i^{(n)}:X(n)\to X(1),
    \qquad
    1\leq i\leq n,
\]
for restriction to the \(i\)-th elementary edge. Write
\[
    \operatorname{seg}_{X,n}:
    X(n)
    \xrightarrow{\cong}
    \underbrace{
    X(1)\times_{X(0)}\cdots\times_{X(0)}X(1)
    }_{n\text{ factors}}
\]
for the edge-restriction isomorphism, and
\[
    \operatorname{gl}_{X,n}:=\operatorname{seg}_{X,n}^{-1}
\]
for its inverse.

For \(X(2)\), abbreviate
\[
    \pi_1:=\pi_1^{(2)},
    \qquad
    \pi_2:=\pi_2^{(2)}.
\]
Let
\[
    \pi_{12}:X(3)\to X(2),
    \qquad
    \pi_{23}:X(3)\to X(2)
\]
be restriction to the first two and last two edges.

\begin{definition}[Segal composition structure]
\label{def:segal-enhancement}
A \textbf{Segal composition structure} on a trajectory sheaf \(X\) consists
of maps
\[
    e_X:X(0)\to X(1),
    \qquad
    c_X:X(2)\to X(1)
\]
satisfying the following equations.

Identities have the correct endpoints:
\[
    s_Xe_X=1_{X(0)},
    \qquad
    t_Xe_X=1_{X(0)}.
\]
Contraction has the endpoints of the composite path:
\[
    s_Xc_X=s_X\pi_1,
    \qquad
    t_Xc_X=t_X\pi_2.
\]
Define
\[
    \ell_X
    :=
    \operatorname{gl}_{X,2}
    \left\langle
        e_Xs_X,
        1_{X(1)}
    \right\rangle,
\]
\[
    r_X
    :=
    \operatorname{gl}_{X,2}
    \left\langle
        1_{X(1)},
        e_Xt_X
    \right\rangle.
\]
The unit equations are
\[
    c_X\ell_X=1_{X(1)},
    \qquad
    c_Xr_X=1_{X(1)}.
\]
Finally, the associativity equation is
\[
    c_X\operatorname{gl}_{X,2}
    \left\langle
        c_X\pi_{12},
        \pi_3^{(3)}
    \right\rangle
    =
    c_X\operatorname{gl}_{X,2}
    \left\langle
        \pi_1^{(3)},
        c_X\pi_{23}
    \right\rangle.
\]
\end{definition}

The endpoint equations ensure that every displayed pairing lands in the
required pullback. The unit equations contract an identity followed by an
arrow and an arrow followed by an identity. The associativity equation says
that the two ways of contracting a three-edge path agree.

\begin{theorem}[Classification of Segal composition structures]
\label{thm:classification-segal-enhancements}
For a trajectory sheaf \(X\), the following data are canonically equivalent:
\begin{enumerate}
    \item Segal composition structures on \(X\);
    \item internal category structures on the reconstructed graph \(G_X\);
    \item isomorphism classes of pairs \((Y,\theta)\), where \(Y\) is an
    augmented strict Segal object and
    \[
        \theta:j^\ast Y\xrightarrow{\cong}X
    \]
    is a specified isomorphism, with isomorphisms taken over \(X\).
\end{enumerate}
\end{theorem}

\begin{proof}
Under the canonical sheaf isomorphism
\[
    \operatorname{seg}_{X,2}:
    X(2)
    \xrightarrow{\cong}
    X(1)\times_{X(0)}X(1),
\]
define
\[
    m_X
    :=
    c_X\operatorname{gl}_{X,2}:
    X(1)\times_{X(0)}X(1)
    \to X(1).
\]
The endpoint equations for \(c_X\) are exactly the source and target axioms
for \(m_X\). The maps \(\ell_X\) and \(r_X\) are the composable pairs formed
by an identity followed by an arrow and by an arrow followed by an identity.
Their contraction equations are the left- and right-unit axioms. The
associativity equation states equality of the two parenthesized composites of
a three-arrow path. Hence a Segal composition structure is precisely an
internal category structure on \(G_X\).

We next construct the corresponding augmented simplicial object explicitly.
Set
\[
    Y_n:=X(n)
\]
for \(n\geq0\), and set
\[
    Y_{-1}:=1.
\]
Define iterated contraction maps
\[
    c_0:=e_X:X(0)\to X(1),
    \qquad
    c_1:=1_{X(1)}.
\]
For \(r\geq1\), let
\[
    \rho_{0,r}:X(r+1)\to X(r)
\]
be restriction to the first \(r\) edges, and define recursively
\[
    c_{r+1}
    :=
    c_X\operatorname{gl}_{X,2}
    \left\langle
        c_r\rho_{0,r},
        \pi_{r+1}^{(r+1)}
    \right\rangle.
\]
The pairing is well typed because the target of the contracted initial block
is the source of the final edge. Associativity shows that \(c_r\) is
independent of parenthesization, and the unit equations handle blocks of
length zero.

Let
\[
    \varphi:[m]\to[n]
\]
be order preserving, with \(m,n\geq0\). An \(n\)-path has vertices
\[
    x_0,x_1,\ldots,x_n.
\]
Define its image under \(Y(\varphi)\) to have vertices
\[
    x_{\varphi(0)},
    x_{\varphi(1)},
    \ldots,
    x_{\varphi(m)}.
\]
For \(1\leq j\leq m\), the \(j\)-th arrow is obtained as follows. If
\[
    \varphi(j-1)<\varphi(j),
\]
restrict to the contiguous block of edges between these two vertices and
apply the contraction
\[
    c_{\varphi(j)-\varphi(j-1)}.
\]
If
\[
    \varphi(j-1)=\varphi(j),
\]
restrict to the repeated vertex and apply \(c_0=e_X\), thereby inserting an
identity. The endpoint equations make the resulting arrows composable, and
the inverse Segal isomorphism
\[
    \operatorname{gl}_{X,m}
\]
assembles them into a morphism
\[
    Y(\varphi):X(n)\to X(m).
\]

For the unique map
\[
    u_n:[-1]\to[n]
\]
in \(\Delta_+\), define
\[
    Y(u_n):Y_n\to Y_{-1}
\]
to be the unique map to \(1\). The identity of \([-1]\) is sent to
\(1_1\).

Suppose
\[
    [\ell]\xrightarrow{\beta}[m]
    \xrightarrow{\varphi}[n].
\]
The blocks determined by \(\varphi\beta\) are obtained by grouping together
consecutive blocks determined first by \(\varphi\) and then by \(\beta\).
Associativity says that contracting the smaller blocks and then contracting
the resulting larger block agrees with contracting the full block directly.
The unit equations remove all zero-length blocks. Therefore
\[
    Y(\varphi\beta)
    =
    Y(\beta)Y(\varphi),
\]
and the construction defines a functor
\[
    Y:\Delta_+^{\mathrm{op}}\to\mathcal E.
\]

If \(\varphi\) is a contiguous interval inclusion of positive length, every
selected block has length one, so \(Y(\varphi)\) is exactly the original
contiguous restriction map of \(X\). Point inclusions select the corresponding
vertex, and maps involving \([-1]\) are uniquely determined. Hence
\[
    j^\ast Y\cong X.
\]
The Segal maps of \(Y\) are precisely the sheaf isomorphisms
\[
    X(n)
    \cong
    X(1)\times_{X(0)}\cdots\times_{X(0)}X(1),
\]
and \(Y_{-1}=1\), so \(Y\) is augmented strict Segal.

Conversely, let \(Y\) be augmented strict Segal and let
\[
    \theta:j^\ast Y\xrightarrow{\cong}X.
\]
Transport the codegeneracy
\[
    Y(\sigma^0):Y_0\to Y_1
\]
and the inner-face map
\[
    Y(\delta^1):Y_2\to Y_1
\]
across \(\theta\). These give \(e_X\) and \(c_X\). The simplicial identities
give the endpoint, unit, and associativity equations. The two constructions
are inverse by Theorem~\ref{thm:internal-nerve-equivalence}.

Finally, an isomorphism of pairs over \(X\) is determined in degrees zero and
one. The Segal isomorphisms determine it in every nonnegative degree, while
the degree \(-1\) component is unique because that object is terminal. Thus
passing to isomorphism classes loses no additional structure.
\end{proof}

\begin{corollary}[Canonical Segal composition of executions]
\label{cor:segal-enhancement-internal-category}
Every internal category \(\mathbb C\), and in particular every execution
category \(\operatorname{Exec}(\Phi)\), determines
\begin{enumerate}
    \item an underlying trajectory sheaf \(\operatorname{Path}(\mathbb C)\);
    \item a canonical Segal composition structure supplied by
    \(N_+\mathbb C\).
\end{enumerate}
\end{corollary}

\subsection{Finite-limit preservation}
\label{subsec:path-nerve-finite-limits}

An internal functor
\[
    F:\mathbb C\to\mathbb D
\]
induces a simplicial map
\[
    N_+F:N_+\mathbb C\to N_+\mathbb D
\]
and, after restriction along \(j\), a morphism of trajectory sheaves
\[
    \operatorname{Path}(F):
    \operatorname{Path}(\mathbb C)
    \to
    \operatorname{Path}(\mathbb D).
\]

\begin{proposition}[Finite-limit preservation]
\label{prop:temp-finite-limit-preservation}
The functors
\[
    N_+:
    \mathsf{IntCat}(\mathcal E)
    \to
    [\Delta_+^{\mathrm{op}},\mathcal E]
\]
and
\[
    \operatorname{Path}:
    \mathsf{IntCat}(\mathcal E)
    \to
    \mathbf{TrajSh}_{\mathcal E}
\]
preserve finite limits.
\end{proposition}

\begin{proof}
Finite limits of internal categories are computed on object and arrow objects,
with identities and composition induced by the relevant universal properties.
Every nerve degree is obtained from those objects by a finite iterated
pullback. Therefore the nerve preserves finite limits degreewise.

Limits of \(\mathcal E\)-valued sheaves are computed pointwise in the
presheaf category. Restriction along \(j\) preserves all pointwise limits.
Hence the underlying trajectory-sheaf functor preserves finite limits.
\end{proof}

\begin{remark}[Segal contraction and polynomial multiplicativity]
\label{rem:segal-to-polynomial-multiplicativity}
The following chapter curries labelled one-step execution data by dependent
product.  Under that transposition, the categorical laws isolated here become
the multiplicativity laws of the resulting response coalgebra:
\[
\begin{array}{c|c}
    \text{Segal execution structure}
    &
    \text{curried response structure}
    \\
    \hline
    \text{identity insertion}
    &
    \text{unit response law}
    \\
    \text{contraction of adjacent executions}
    &
    \text{multiplicative response law}
    \\
    \text{unit and associativity of contraction}
    &
    \text{multiplicativity coherence}.
\end{array}
\]
This is a change of presentation rather than a second execution semantics.
The polynomial coalgebra records total one-step responses in curried form;
flattening it recovers the execution category whose Segal structure is
described in this section.
\end{remark}

\section{Labelled finite trajectories}
\label{sec:temporal-execution-semantics}

The finite execution objects constructed in the opening sections are now
recognized as the underlying trajectory sheaf of the execution category. The
abstract equivalences therefore return directly to the automata-theoretic
setting.

\subsection{The execution trajectory sheaf}
\label{subsec:temporal-execution-sheaf}

\begin{definition}[Execution trajectory sheaf]
\label{def:temporal-execution-sheaf}
The \textbf{execution trajectory sheaf} of
\[
    \Phi:\mathbb A\rightsquigarrow\mathbb B
\]
is
\[
    \mathscr T_\Phi
    :=
    \operatorname{Path}
    \bigl(
        \operatorname{Exec}(\Phi)
    \bigr).
\]
\end{definition}

Thus
\[
    \mathscr T_\Phi(-1)=1,
    \qquad
    \mathscr T_\Phi(0)=P,
\]
\[
    \mathscr T_\Phi(1)
    =
    A_1\times_{t_A,A_0,p_P}P,
\]
and for \(n\geq2\),
\[
    \mathscr T_\Phi(n)
    =
    \underbrace{
    \mathscr T_\Phi(1)
    \times_P\cdots\times_P
    \mathscr T_\Phi(1)
    }_{n\text{ factors}}.
\]
These are exactly the finite execution objects already derived from the Mealy
transition. The sheaf structure says that an execution is recovered uniquely
from its one-step segments and their common state boundaries.

The input and output-comparison labelling functors induce morphisms of
trajectory sheaves
\[
    \operatorname{Path}(I_\Phi):
    \mathscr T_\Phi
    \to
    \operatorname{Path}(\mathbb A^{\mathrm{op}}),
\]
\[
    \operatorname{Path}(O_\Phi):
    \mathscr T_\Phi
    \to
    \operatorname{Path}(\mathbb B^{\mathrm{op}}).
\]

\begin{definition}[Labelled finite-trajectory semantics]
\label{def:labelled-temporal-execution-semantics}
The \textbf{labelled finite-trajectory semantics} of \(\Phi\) is the span
\[
\begin{tikzcd}[column sep=large]
    \operatorname{Path}(\mathbb A^{\mathrm{op}})
    &
    \mathscr T_\Phi
        \arrow[l, "{\operatorname{Path}(I_\Phi)}"']
        \arrow[r, "{\operatorname{Path}(O_\Phi)}"]
    &
    \operatorname{Path}(\mathbb B^{\mathrm{op}}).
\end{tikzcd}
\]
\end{definition}

At path length \(n\), the left map records the \(n\)-step input path, the
right map records the \(n\)-step output-comparison path, and the middle object
records the intervening state trajectory.

\subsection{Deterministic path lifting in sheaf form}
\label{subsec:deterministic-path-lifting-sheaf-form}

For \(n\geq0\), let
\[
    \lambda_0^{A,n}:
    \operatorname{Path}(\mathbb A^{\mathrm{op}})(n)
    \to A_0
\]
denote restriction to the initial boundary of the input path.
Proposition~\ref{prop:deterministic-path-lifting} may now be written
canonically as the pullback isomorphism
\[
\begin{tikzcd}[column sep=large]
    \mathscr T_\Phi(n)
        \arrow[r, "\cong"]
    &
    P\times_{A_0}
    \operatorname{Path}(\mathbb A^{\mathrm{op}})(n).
\end{tikzcd}
\]
Equivalently, the square
\[
\begin{tikzcd}
    \mathscr T_\Phi(n)
        \arrow[r, "{\operatorname{Path}(I_\Phi)_n}"]
        \arrow[d, "s_n^\Phi"']
    &
    \operatorname{Path}(\mathbb A^{\mathrm{op}})(n)
        \arrow[d, "{\lambda_0^{A,n}}"]
    \\
    P
        \arrow[r, "p_P"']
    &
    A_0
\end{tikzcd}
\]
is a pullback.

This square is the finite-run determinism property of the trajectory
semantics. It says that the middle object of the labelled span contains no
choice beyond the initial state once the input path is fixed.

\subsection{Segal-complete labelled execution}
\label{subsec:segal-enhanced-execution-semantics}

The trajectory span forgets no finite runs, but it retains only contiguous
restriction. The execution categories also provide identity insertion and
contraction. These operations are recorded by the augmented nerves.

\begin{definition}[Segal-complete labelled execution]
\label{def:segal-enhanced-execution-semantics}
The \textbf{Segal-complete labelled execution semantics} of \(\Phi\) is the
span of augmented strict Segal objects
\[
\begin{tikzcd}[column sep=large]
    N_+(\mathbb A^{\mathrm{op}})
    &
    N_+\operatorname{Exec}(\Phi)
        \arrow[l, "N_+I_\Phi"']
        \arrow[r, "N_+O_\Phi"]
    &
    N_+(\mathbb B^{\mathrm{op}}).
\end{tikzcd}
\]
\end{definition}

Restriction along
\[
    j:\mathbf{Int}_{\mathrm{disc}}\hookrightarrow\Delta_+
\]
recovers the labelled finite-trajectory span. The trajectory span records
finite runs, contiguous restriction, truncation, and gluing. The
Segal-complete span additionally records identity executions, contraction of
adjacent steps, and coherence of iterated contraction.

\subsection{Functoriality in strict semantic homomorphisms}
\label{subsec:trajectory-functoriality-strict}

For a strict semantic homomorphism
\[
    f:\Phi\to\Psi,
\]
the internal functor \(\operatorname{Exec}(f)\) induces
\[
    \mathscr T_f:
    \mathscr T_\Phi\to\mathscr T_\Psi
\]
and
\[
    N_+\operatorname{Exec}(f):
    N_+\operatorname{Exec}(\Phi)
    \to
    N_+\operatorname{Exec}(\Psi).
\]
Both commute with input and output-comparison labels.

At length \(n\), the map \(\mathscr T_f(n)\) sends
\[
\begin{tikzcd}[column sep=large]
    x_0
        \arrow[r]
    &
    x_1
        \arrow[r]
    &
    \cdots
        \arrow[r]
    &
    x_n
\end{tikzcd}
\]
to
\[
\begin{tikzcd}[column sep=large]
    f(x_0)
        \arrow[r]
    &
    f(x_1)
        \arrow[r]
    &
    \cdots
        \arrow[r]
    &
    f(x_n),
\end{tikzcd}
\]
with exactly the same input and output-comparison paths.

\subsection{Composition of trajectory spans}
\label{subsec:temporal-composition}

Applying the finite-limit-preserving functor \(\operatorname{Path}\) to
Theorem~\ref{thm:execution-category-composition} gives the composition law for
finite trajectories.

\begin{theorem}[Composition of labelled finite trajectories]
\label{thm:temporal-composition}
For composable Mealy morphisms
\[
    \Phi:\mathbb A\rightsquigarrow\mathbb B,
    \qquad
    \Psi:\mathbb B\rightsquigarrow\mathbb C,
\]
there is a canonical isomorphism of trajectory sheaves
\[
\begin{tikzcd}[column sep=large]
    \mathscr T_{\Psi\circ\Phi}
        \arrow[r, "\cong"]
    &
    \mathscr T_\Phi
        \times_{\operatorname{Path}(\mathbb B^{\mathrm{op}})}
    \mathscr T_\Psi.
\end{tikzcd}
\]
Equivalently, for every \(n\geq-1\),
\[
    \mathscr T_{\Psi\circ\Phi}(n)
    \cong
    \mathscr T_\Phi(n)
    \times_{\operatorname{Path}(\mathbb B^{\mathrm{op}})(n)}
    \mathscr T_\Psi(n).
\]
\end{theorem}

\begin{proof}
The execution category of the composite is the pullback of the two execution
categories over \(\mathbb B^{\mathrm{op}}\). The functor
\(\operatorname{Path}\) preserves finite limits, so it sends this pullback
to the displayed pullback of trajectory sheaves. Limits of sheaves are
computed pointwise, giving the levelwise formula.
\end{proof}

Thus an \(n\)-step execution of the composite consists precisely of
\begin{enumerate}
    \item an \(n\)-step execution of \(\Phi\);
    \item an \(n\)-step execution of \(\Psi\);
    \item equality between the output-comparison path of \(\Phi\) and the
    input path of \(\Psi\).
\end{enumerate}
The equality is pathwise, not merely equality of the composites of those
paths. The chosen segmentations of the two runs must agree at every step.

\begin{proposition}[Naturality of trajectory composition]
\label{prop:naturality-temporal-composition}
The isomorphism of
Theorem~\ref{thm:temporal-composition}
is natural in strict semantic homomorphisms of both Mealy factors.
\end{proposition}

\begin{proof}
Naturality was established at the level of execution categories in
Proposition~\ref{prop:naturality-execution-composition}. Applying the functor
\(\operatorname{Path}\) to the commuting comparison square gives the
required square of trajectory sheaves.
\end{proof}

The identity Mealy morphism has trajectory semantics equal to the identity
span on \(\operatorname{Path}(\mathbb A^{\mathrm{op}})\). The associativity
and unit comparisons are inherited from the canonical pullback comparisons.
Consequently, labelled finite-trajectory semantics defines a pseudofunctor to
the bicategory of spans in \(\mathbf{TrajSh}_{\mathcal E}\).

\section{Quotients and trajectory minimization}
\label{sec:temporal-quotients-minimization}

\subsection{Strict regular quotients}
\label{subsec:strict-regular-quotients-temporal}

Let
\[
    e:\Phi\to\Psi
\]
be a strict semantic homomorphism between Mealy morphisms over the same input
and output interfaces.

\begin{definition}[Strict regular quotient]
\label{def:strict-regular-quotient-temporal}
The homomorphism \(e\) is a \textbf{strict regular quotient} if its carrier
map
\[
    e:P\twoheadrightarrow Q
\]
is a regular epimorphism in
\[
    \mathcal E/(A_0\times B_0).
\]
Equivalently, its underlying map in \(\mathcal E\) is a regular epimorphism.
\end{definition}

\begin{proposition}[Execution preserves strict regular quotients]
\label{prop:execution-preserves-regular-quotients}
If
\[
    e:\Phi\twoheadrightarrow\Psi
\]
is a strict regular quotient, then
\[
    \operatorname{Exec}(e):
    \operatorname{Exec}(\Phi)
    \to
    \operatorname{Exec}(\Psi)
\]
is regular epic on both object and arrow objects.
\end{proposition}

\begin{proof}
On objects, the map is \(e:P\twoheadrightarrow Q\). On arrows, it is
\[
\begin{tikzcd}[column sep=large]
    A_1\times_{A_0}P
        \arrow[r, "{1_{A_1}\times e}"]
    &
    A_1\times_{A_0}Q.
\end{tikzcd}
\]
Because \(e\) is a map over \(A_0\), this arrow map is a pullback of \(e\).
Regular epimorphisms are stable under pullback in a Grothendieck topos, so the
arrow map is regular epic.
\end{proof}

\subsection{Regular quotients in every finite path length}
\label{subsec:levelwise-trajectory-quotients}

\begin{theorem}[Trajectory preservation of regular quotients]
\label{thm:temporal-preservation-regular-quotients}
Let
\[
    e:\Phi\twoheadrightarrow\Psi
\]
be a strict regular quotient. Then, for every \(n\geq-1\),
\[
    \mathscr T_e(n):
    \mathscr T_\Phi(n)
    \twoheadrightarrow
    \mathscr T_\Psi(n)
\]
is a regular epimorphism.
\end{theorem}

\begin{proof}
At length \(-1\), the map is the identity of \(1\). For \(n\geq0\),
deterministic path lifting gives the two pullback descriptions
\[
    \mathscr T_\Phi(n)
    \cong
    P\times_{A_0}\operatorname{Path}(\mathbb A^{\mathrm{op}})(n),
\]
\[
    \mathscr T_\Psi(n)
    \cong
    Q\times_{A_0}\operatorname{Path}(\mathbb A^{\mathrm{op}})(n).
\]
Under these isomorphisms, \(\mathscr T_e(n)\) is
\[
\begin{tikzcd}[column sep=large]
    P\times_{A_0}\operatorname{Path}(\mathbb A^{\mathrm{op}})(n)
        \arrow[r, "{e\times_{A_0}1}"]
    &
    Q\times_{A_0}\operatorname{Path}(\mathbb A^{\mathrm{op}})(n).
\end{tikzcd}
\]
This map is a pullback of \(e\), hence is a regular epimorphism.
\end{proof}

\begin{corollary}[Levelwise quotient of trajectory sheaves]
\label{cor:levelwise-quotient-trajectory-sheaves}
A strict regular quotient of Mealy systems induces a levelwise regular
epimorphism of trajectory sheaves
\[
    \mathscr T_\Phi
    \twoheadrightarrow
    \mathscr T_\Psi.
\]
\end{corollary}

\subsection{Behavioural minimization of finite trajectories}
\label{subsec:temporal-semantic-minimization}

Recall the semantic image factorization
\[
\begin{tikzcd}[column sep=large]
    P
        \arrow[r, two heads, "\eta_\Phi"]
    &
    \operatorname{Beh}(\Phi)
        \arrow[r, tail]
    &
    \mathbf{Beh}^{\mathrm{can}}_{\mathbb A,\mathbb B,0}.
\end{tikzcd}
\]
The carrier \(\operatorname{Beh}(\Phi)\) inherits a canonical strict Mealy
structure, and
\[
    \eta_\Phi:
    \Phi\twoheadrightarrow\operatorname{Beh}(\Phi)
\]
is a strict regular quotient.

\begin{theorem}[Behavioural quotient of finite trajectories]
\label{thm:temporal-semantic-quotient}
Every Mealy morphism \(\Phi\) admits a canonical levelwise regular quotient
\[
\begin{tikzcd}[column sep=large]
    \mathscr T_\Phi
        \arrow[r, two heads]
    &
    \mathscr T_{\operatorname{Beh}(\Phi)}.
\end{tikzcd}
\]
At path length \(n\geq0\), two executions have the same image precisely when
they lie over the same input path and their initial states have the same
residual semantic image. Equivalently, along a common input path, their state
images agree at every boundary.
\end{theorem}

\begin{proof}
Apply
Theorem~\ref{thm:temporal-preservation-regular-quotients}
to \(\eta_\Phi\).

For the kernel description, deterministic path lifting identifies the
length-\(n\) quotient with
\[
    \eta_\Phi\times_{A_0}1.
\]
Its kernel pair therefore requires equality of the input-path component and
equality of the semantic images of the initial states. Strict transition
preservation forces equality of the semantic images of every subsequent state
along that common input path. Conversely, equality at every boundary includes
equality at the initial boundary and therefore implies equality under the
quotient map.
\end{proof}

\begin{corollary}[Separated trajectory semantics]
\label{cor:separated-temporal-semantics}
If \(\Phi\) is behaviourally separated, then
\[
    \mathscr T_\Phi
    \cong
    \mathscr T_{\operatorname{Beh}(\Phi)}.
\]
\end{corollary}

\begin{proof}
For a behaviourally separated system, the semantic quotient \(\eta_\Phi\) is
an isomorphism. The execution-category and trajectory constructions preserve
isomorphisms.
\end{proof}

Thus behavioural separatedness removes duplicate labelled trajectories that
arise from behaviourally indistinguishable initial states over the same input
path. It does not quotient different segmentations of a path; those remain
related by Segal contraction rather than identified in the trajectory sheaf.

\subsection{Finite-trajectory Myhill--Nerode theorem}
\label{subsec:temporal-myhill-nerode}

Let \(K\) be a specification, and let
\[
    \operatorname{Der}_0(K)
    \hookrightarrow
    \mathbf{Beh}^{\mathrm{can}}_{\mathbb A,\mathbb B,0}
\]
be its derivative-generated residual theory. Since
\(\operatorname{Der}_0(K)\) is derivative invariant, the canonical
transition and output-comparison operations restrict to it. Write
\[
    \iota\bigl(\operatorname{Der}_0(K)\bigr)
\]
for the resulting canonical Mealy realization.

Suppose \(\Phi\) realizes \(K\) and is reachable in the sense that there is a
canonical isomorphism
\[
    \operatorname{Beh}(\Phi)
    \cong
    \iota\bigl(\operatorname{Der}_0(K)\bigr).
\]

\begin{corollary}[Finite-trajectory Myhill--Nerode theorem]
\label{cor:temporal-reachable-quotient}
Every reachable realization of \(K\) admits a canonical levelwise regular
quotient
\[
\begin{tikzcd}[column sep=large]
    \mathscr T_\Phi
        \arrow[r, two heads]
    &
    \mathscr T_{\iota(\operatorname{Der}_0(K))}.
\end{tikzcd}
\]
If \(\Phi\) is also behaviourally separated, then this map is an
isomorphism.
\end{corollary}

\begin{proof}
Substitute the reachable identification into
Theorem~\ref{thm:temporal-semantic-quotient}. In the separated case, apply
Corollary~\ref{cor:separated-temporal-semantics}.
\end{proof}

This is the finite-trajectory form of categorical Myhill--Nerode theory:
reachable and separated realizations have exactly the labelled trajectories of
the canonical derivative realization.

\section{Residual trajectories and endpoint reachability}
\label{sec:residual-contracts-temporal-safety}

The preceding section applied state minimization to arbitrary Mealy
morphisms.  We now specialize to the canonical behaviour machine and
interpret derivative invariance and generated derivative closure directly in
terms of finite residual executions.  The universal predicate-transformer
form of finite-path safety is postponed to the later dynamic-logic chapter,
where it can be compared with span boxes and polynomial profile boxes.

\subsection{Canonical residual trajectories}
\label{subsec:canonical-residual-trajectories}

Fix input and output categories
\[
    \mathbb A,
    \qquad
    \mathbb B.
\]
Let
\[
    \mathbf{Beh}^{\mathrm{can}}_{\mathbb A,\mathbb B}
\]
be the canonical behaviour Mealy morphism, with carrier
\[
    \mathbf{Beh}_0
    :=
    \mathbf{Beh}_{\mathbb A,\mathbb B,0}.
\]
Its execution trajectory sheaf is
\[
    \mathscr T_{\mathbf{Beh}}
    :=
    \mathscr T_{\mathbf{Beh}^{\mathrm{can}}_{\mathbb A,\mathbb B}}.
\]

For \(n\geq0\) and \(0\leq i\leq n\), write
\[
    \nu_i^{(n)}:
    \mathscr T_{\mathbf{Beh}}(n)
    \to
    \mathbf{Beh}_0
\]
for the \(i\)-th residual-state map. In particular,
\[
    s_n:=\nu_0^{(n)},
    \qquad
    t_n:=\nu_n^{(n)}.
\]
A generalized element has the form
\[
\begin{tikzcd}[column sep=large]
    H_0
        \arrow[r, "a_1"]
    &
    H_1
        \arrow[r, "a_2"]
    &
    \cdots
        \arrow[r, "a_n"]
    &
    H_n,
\end{tikzcd}
\]
where
\[
    H_i=\partial_{a_i}H_{i-1}.
\]
The path is labelled in \(\mathbb A^{\mathrm{op}}\) by the opposites of the
input arrows and in \(\mathbb B^{\mathrm{op}}\) by the corresponding output
comparisons.

\begin{remark}[Ambient subobject lattices in the trajectory formulas]
\label{rem:temporal-subobject-lattices}
Behavioural predicates were introduced as subobjects of \(\mathbf{Beh}_0\) in
the slice
\[
    \mathcal E/(A_0\times B_0).
\]
The trajectory state maps
\[
    s_n,\nu_i^{(n)}:
    \mathscr T_{\mathbf{Beh}}(n)
    \to
    \mathbf{Beh}_0
\]
are not generally morphisms in that fixed slice, because the input and output
anchors may change along an execution.

Accordingly, the pullbacks, regular images, and dependent products in this
section are computed in \(\mathcal E\). This introduces no ambiguity: for
the fixed structure map
\[
    \mathbf{Beh}_0\to A_0\times B_0,
\]
there is a canonical identification
\[
    \operatorname{Sub}_{\mathcal E}(\mathbf{Beh}_0)
    \cong
    \operatorname{Sub}_{\mathcal E/(A_0\times B_0)}
    (\mathbf{Beh}_0).
\]
Every resulting subobject of \(\mathbf{Beh}_0\) inherits its structure map to
\(A_0\times B_0\).
\end{remark}

\subsection{Derivative-invariant behavioural specifications}
\label{subsec:trajectory-realization-residual-contract}

Let
\[
    V\hookrightarrow\mathbf{Beh}_0
\]
be derivative invariant. The canonical transition and output-comparison maps
restrict to \(V\), producing the canonical realization \(\iota(V)\).

\begin{theorem}[Trajectory object of a derivative-invariant specification]
\label{thm:trajectory-object-residual-contract}
For every \(n\geq0\), there is a canonical pullback isomorphism
\[
\begin{tikzcd}[column sep=large]
    \mathscr T_{\iota(V)}(n)
        \arrow[r, "\cong"]
    &
    V\times_{\mathbf{Beh}_0,s_n}
    \mathscr T_{\mathbf{Beh}}(n).
\end{tikzcd}
\]
Consequently, a canonical residual trajectory lies in
\(\mathscr T_{\iota(V)}(n)\) if and only if its initial residual state lies
in \(V\).
\end{theorem}

\begin{proof}
Deterministic path lifting gives
\[
    \mathscr T_{\mathbf{Beh}}(n)
    \cong
    \mathbf{Beh}_0
    \times_{A_0}
    \operatorname{Path}(\mathbb A^{\mathrm{op}})(n),
\]
under which \(s_n\) is projection to \(\mathbf{Beh}_0\). Similarly,
\[
    \mathscr T_{\iota(V)}(n)
    \cong
    V\times_{A_0}
    \operatorname{Path}(\mathbb A^{\mathrm{op}})(n).
\]
Pulling the first isomorphism back along
\(V\hookrightarrow\mathbf{Beh}_0\) produces the second. Derivative
invariance guarantees that the transition structure restricts and hence that
every later residual state remains in \(V\).
\end{proof}

\begin{corollary}[All trajectory boundaries remain in the specification]
\label{cor:all-vertices-remain-contract}
For every \(n\geq0\) and every \(0\leq i\leq n\),
\[
    s_n^\ast V
    \leq
    \bigl(\nu_i^{(n)}\bigr)^\ast V
\]
as subobjects of \(\mathscr T_{\mathbf{Beh}}(n)\).
\end{corollary}

\begin{proof}
By the theorem, a trajectory beginning in \(V\) is a trajectory of the
restricted canonical realization \(\iota(V)\). Every state projection of
that restricted trajectory factors through \(V\).
\end{proof}

\subsection{Finite-path characterization of derivative invariance}
\label{subsec:temporal-characterization-derivative-invariance}

Let
\[
    C\hookrightarrow\mathbf{Beh}_0
\]
be an arbitrary behavioural predicate.

\begin{theorem}[Finite-path characterization of derivative invariance]
\label{thm:finite-horizon-characterization-contracts}
The following are equivalent:
\begin{enumerate}
    \item \(C\) is derivative invariant;
    \item every one-step canonical residual execution beginning in \(C\)
    ends in \(C\), equivalently
    \[
        s_1^\ast C\leq t_1^\ast C;
    \]
    \item for every \(n\geq0\) and every \(0\leq i\leq n\),
    \[
        s_n^\ast C
        \leq
        \bigl(\nu_i^{(n)}\bigr)^\ast C;
    \]
    \item every finite canonical residual execution beginning in \(C\)
    remains in \(C\) at every state boundary.
\end{enumerate}
\end{theorem}

\begin{proof}
Derivative invariance says exactly that every applicable derivative of a
state in \(C\) lies again in \(C\). The object of applicable one-step
derivatives is \(\mathscr T_{\mathbf{Beh}}(1)\), with source and target maps
\(s_1\) and \(t_1\). Thus conditions one and two are equivalent.

Assume the one-step condition. An \(n\)-step residual execution is a chain of
one-step executions. Repeated application of the one-step inclusion propagates
membership in \(C\) from the initial state to every later boundary. This gives
condition three. Conditions three and four are the external and internal
readings of the same family of subobject inclusions. Finally, condition three
at \(n=1\) and \(i=1\) is the one-step condition.
\end{proof}

Thus derivative-invariant behavioural specifications are exactly the
state-determined predicates closed under every finite residual prefix.

\subsection{Generated derivative closure as endpoint reachability}
\label{subsec:temporal-interpretation-generated-closure}

For a predicate
\[
    C\hookrightarrow\mathbf{Beh}_0,
\]
define the object of length-\(n\) canonical trajectories beginning in \(C\)
by
\[
    \mathscr T_C(n)
    :=
    C\times_{\mathbf{Beh}_0,s_n}
    \mathscr T_{\mathbf{Beh}}(n).
\]
Write
\[
    t_{n,C}:
    \mathscr T_C(n)
    \to
    \mathbf{Beh}_0
\]
for the terminal-state map. Form
\[
    S_C
    :=
    \coprod_{n\in\mathbb N}\mathscr T_C(n)
\]
and
\[
    \tau_C
    :=
    \coprod_{n\in\mathbb N}t_{n,C}:
    S_C\to\mathbf{Beh}_0.
\]

\begin{definition}[Finite-path endpoint reachability]
\label{def:temporal-reachability-closure}
The \textbf{finite-path endpoint reachability closure} of \(C\) is
\[
    \operatorname{Reach}_{\mathrm{path}}(C)
    :=
    \operatorname{Im}(\tau_C).
\]
\end{definition}

The augmented nerve of the canonical behaviour category supplies a total
contraction
\[
    \kappa_n:
    \mathscr T_{\mathbf{Beh}}(n)
    \to
    \mathscr T_{\mathbf{Beh}}(1)
\]
induced by the order-preserving map
\[
    [1]\to[n],
    \qquad
    0\mapsto0,
    \quad
    1\mapsto n.
\]
For \(n=0\), this is identity insertion. It satisfies
\[
    s_1\kappa_n=s_n,
    \qquad
    t_1\kappa_n=t_n.
\]

\begin{theorem}[Generated derivative closure is endpoint reachability]
\label{thm:generated-closure-temporal-reachability}
There is a canonical equality
\[
    \operatorname{Reach}_{\mathrm{path}}(C)
    =
    \operatorname{Gen}_{\partial}(C).
\]
\end{theorem}

\begin{proof}
Because \(s_1\kappa_n=s_n\), total contraction restricts to
\[
    \kappa_{n,C}:
    \mathscr T_C(n)
    \to
    \mathscr T_C(1).
\]
The terminal-state equation gives
\[
    t_{n,C}
    =
    t_{1,C}\kappa_{n,C}.
\]
Hence
\[
    \operatorname{Im}(t_{n,C})
    \leq
    \operatorname{Im}(t_{1,C})
\]
for every \(n\). Conversely, the length-one component occurs among the
coproduct components of \(\tau_C\). Therefore
\[
    \operatorname{Im}(\tau_C)
    =
    \operatorname{Im}(t_{1,C}).
\]

The map \(t_{1,C}\) sends a behaviour \(H\in C\), together with an
applicable input arrow \(a\), to the derivative \(\partial_aH\). By the
one-step construction of generated derivative closure in the preceding
chapter, its regular image is exactly
\[
    \operatorname{Gen}_{\partial}(C).
\]
Therefore
\[
    \operatorname{Reach}_{\mathrm{path}}(C)
    =
    \operatorname{Gen}_{\partial}(C).
\]
\end{proof}

The theorem also explains why all finite lengths add no new endpoint states
beyond length one in the full categorical execution semantics. Any finite
input path has a composite input arrow, and total contraction processes that
composite in one categorical step.

\begin{remark}[Deferred universal safety construction]
\label{rem:residual-universal-safety-deferred}
The inclusions
\[
    s_n^\ast C
    \leq
    \bigl(\nu_i^{(n)}\bigr)^\ast C
\]
already give the finite-trajectory characterization of derivative
invariance.  To construct the largest derivative-invariant subobject below an
arbitrary predicate \(C\), one additionally quantifies universally over all
trajectory lengths and all boundary maps.  That construction uses the
dependent products and complete subobject lattices appropriate to the
logical layer of the theory.  It will therefore be developed with the
trajectory box modalities in the later dynamic-logic chapter rather than
repeated here.
\end{remark}

\section{One-object, Boolean, and generator-level specializations}
\label{sec:temporal-one-object-boolean}

\subsection{One-object input and output categories}
\label{subsec:temporal-one-object-specialization}

Suppose
\[
    \mathbb A=\mathbf B M,
    \qquad
    \mathbb B=\mathbf B N
\]
are one-object internal categories associated with monoid objects \(M\) and
\(N\). Using the canonical symmetry
\[
    M\times P\cong P\times M,
\]
we write execution arrows in state-first form
\[
    (x,m).
\]

A Mealy morphism consists of a state object \(P\), a right \(M\)-action
\[
    x\cdot m,
\]
and an output cocycle
\[
    \omega:M\times P\to N
\]
satisfying
\[
    x\cdot(mn)
    =
    (x\cdot m)\cdot n
\]
and
\[
    \omega(mn,x)
    =
    \omega(m,x)\,
    \omega(n,x\cdot m).
\]

The execution category has object object \(P\) and arrow object
\(P\times M\). Its source, target, identity, and composition are
\[
    s(x,m)=x,
    \qquad
    t(x,m)=x\cdot m,
\]
\[
    e(x)=(x,1_M),
\]
\[
    (x,m);(x\cdot m,n)
    =
    (x,mn).
\]
The input label of \((x,m)\) is \(m^{\mathrm{op}}\), and the output label is
\(\omega(m,x)^{\mathrm{op}}\).

An \(n\)-step execution is uniquely determined by
\[
    (x;m_1,\ldots,m_n).
\]
Its state path is
\[
\begin{tikzcd}[column sep=large]
    x
        \arrow[r, "m_1"]
    &
    x\cdot m_1
        \arrow[r, "m_2"]
    &
    \cdots
        \arrow[r, "m_n"]
    &
    x\cdot(m_1\cdots m_n).
\end{tikzcd}
\]
The output-comparison arrows in \(N\) are
\[
    \omega(m_1,x),
\]
\[
    \omega(m_2,x\cdot m_1),
\]
\[
    \ldots,
\]
\[
    \omega
    \bigl(
        m_n,
        x\cdot(m_1\cdots m_{n-1})
    \bigr).
\]
The output path in \(\mathbf BN^{\mathrm{op}}\) consists of their opposites.

Contiguous restriction selects a consecutive block of the monoid labels and
the corresponding boundary state. Gluing concatenates two label strings when
the terminal state of the first equals the initial state of the second.
Contraction sends adjacent labels \(m_i,m_{i+1}\) to \(m_im_{i+1}\). The
cocycle law says that the output comparison of the contracted execution is the
composite of the two original output comparisons.

Total contraction sends
\[
    (x;m_1,\ldots,m_n)
\]
to
\[
    (x;m_1\cdots m_n).
\]
This concrete formula exhibits directly why path length is segmentation in the
full one-object category: an arbitrary word or monoid element already labels a
single categorical step.

\subsection{Boolean derivative executions}
\label{subsec:boolean-language-trajectories}

From this subsection onward, work in
\[
    \mathcal E=\mathbf{Set},
\]
and let \(\Sigma\) be an ordinary set of letters. Let
\[
    L:\Sigma^\ast\to2
\]
be a Boolean language. Write
\[
    \operatorname{Der}(L)
    :=
    \{
        \partial_uL
        \mid
        u\in\Sigma^\ast
    \}
\]
for its set of left derivatives, where
\[
    \partial_uL(w):=L(uw).
\]
The right \(\Sigma^\ast\)-action on residual languages is
\[
    H\cdot v:=\partial_vH.
\]
For \(H=\partial_uL\), one has
\[
    H\cdot v
    =
    \partial_v(\partial_uL)
    =
    \partial_{uv}L.
\]
The observation map is
\[
    o:\operatorname{Der}(L)\to2,
    \qquad
    o(H):=H(\varepsilon).
\]

Because the input category is
\[
    \mathbf B\Sigma^\ast,
\]
a one-step categorical execution may be labelled by an arbitrary word
\(v\in\Sigma^\ast\), not only by a single letter. The full derivative
execution category \(\operatorname{Exec}(L)\) has
\begin{enumerate}
    \item objects \(H\in\operatorname{Der}(L)\);
    \item arrows
    \[
        (H,v):
        H\to\partial_vH,
        \qquad
        v\in\Sigma^\ast.
    \]
\end{enumerate}
Composition is
\[
\begin{tikzcd}[column sep=large]
    H
        \arrow[r, "v"]
    &
    \partial_vH
        \arrow[r, "w"]
    &
    \partial_{vw}H,
\end{tikzcd}
\]
and contracts to the arrow \((H,vw)\).

Write
\[
    \mathscr T_L
    :=
    \operatorname{Path}
    \bigl(
        \operatorname{Exec}(L)
    \bigr)
\]
for the trajectory sheaf of the full word action.

\begin{proposition}[Finite word-labelled derivative executions]
\label{prop:temporal-derivative-automaton}
The \(n\)-sections of \(\mathscr T_L\) are tuples
\[
    (H;v_1,\ldots,v_n),
\]
where
\[
    H\in\operatorname{Der}(L),
    \qquad
    v_i\in\Sigma^\ast.
\]
The associated residual path is
\[
    H,
    \quad
    \partial_{v_1}H,
    \quad
    \partial_{v_2}\partial_{v_1}H,
    \quad
    \ldots,
    \quad
    \partial_{v_n}\cdots\partial_{v_1}H.
\]
If \(H=\partial_uL\), this is
\[
    \partial_uL,
    \quad
    \partial_{uv_1}L,
    \quad
    \partial_{uv_1v_2}L,
    \quad
    \ldots,
    \quad
    \partial_{uv_1\cdots v_n}L.
\]
\end{proposition}

\begin{proof}
Apply deterministic path lifting to the right
\(\Sigma^\ast\)-action on \(\operatorname{Der}(L)\). An input path of
length \(n\) in the one-object category
\(\mathbf B(\Sigma^\ast)^{\mathrm{op}}\) is a tuple
\[
    (v_1,\ldots,v_n)\in(\Sigma^\ast)^n.
\]
The successive states are obtained by the displayed right action. The
calculation
\[
    \partial_{v_i}\partial_{uv_1\cdots v_{i-1}}L
    =
    \partial_{uv_1\cdots v_i}L
\]
gives the second description.
\end{proof}

The Boolean output comparison for a word step may be read as the observation
change along that derivative transition. For the purposes of path structure,
what matters is that the output labels compose according to the Mealy cocycle,
exactly as in the general one-object case.

\subsection{Letterwise execution paths}
\label{subsec:letterwise-derivative-trajectories}

The categorical execution above treats every word as a one-step label. To make
path length count letters, retain only transitions labelled by elements of the
generating set \(\Sigma\).

\begin{definition}[Letterwise derivative graph]
\label{def:letterwise-derivative-trajectory-sheaf}
The \textbf{letterwise derivative graph}
\[
    G_L^\Sigma
\]
has
\begin{enumerate}
    \item vertex set \(\operatorname{Der}(L)\);
    \item edge set \(\operatorname{Der}(L)\times\Sigma\);
    \item source and target maps
    \[
        s(H,a)=H,
        \qquad
        t(H,a)=\partial_aH.
    \]
\end{enumerate}
Its trajectory sheaf is denoted
\[
    \mathscr T_L^\Sigma
    :=
    \operatorname{Path}(G_L^\Sigma).
\]
\end{definition}

An \(n\)-section is
\[
    (H;a_1,\ldots,a_n),
    \qquad
    a_i\in\Sigma.
\]
Starting at \(L\), the residual path is
\[
\begin{tikzcd}[column sep=large]
    L
        \arrow[r, "a_1"]
    &
    \partial_{a_1}L
        \arrow[r, "a_2"]
    &
    \cdots
        \arrow[r, "a_n"]
    &
    \partial_{a_1\cdots a_n}L.
\end{tikzcd}
\]
The observation at boundary \(i\) is
\[
    o\bigl(\partial_{a_1\cdots a_i}L\bigr)
    =
    L(a_1\cdots a_i).
\]
Hence a letterwise execution beginning at \(L\) records the acceptance value
of every prefix of the observed word, including the empty prefix.

\begin{proposition}[Letterwise restriction]
\label{prop:letterwise-derivative-trajectories}
For every \(n\geq0\), there is a canonical monomorphism
\[
    \mathscr T_L^\Sigma(n)
    \hookrightarrow
    \mathscr T_L(n)
\]
whose image consists of the word-labelled executions
\[
    (H;v_1,\ldots,v_n)
\]
for which every \(v_i\) has length one. These monomorphisms commute with
contiguous path restriction.
\end{proposition}

\begin{proof}
The inclusion of generators
\[
    \Sigma\hookrightarrow\Sigma^\ast
\]
induces an inclusion of the letter-transition graph into the underlying graph
of the full \(\Sigma^\ast\)-action category. Applying the graph path functor
gives the levelwise monomorphisms. Contiguous restriction preserves the
property that every one-step label has length one.
\end{proof}

The letterwise trajectory sheaf is closed under restriction and gluing. It is
not closed under the contraction inherited from the full execution category:
contracting two letter-labelled edges produces an edge labelled by a word of
length two, which is not an edge of \(G_L^\Sigma\).

\subsection{Free categorical completion of letterwise execution}
\label{subsec:letterwise-free-category-completion}

Let
\[
    \operatorname{Free}(G_L^\Sigma)
\]
be the free category on the letterwise derivative graph. An arrow is a finite
letter path
\[
    (H;a_1,\ldots,a_n),
\]
including the empty path at \(H\).

\begin{theorem}[Word-labelled execution as free categorical completion]
\label{thm:letterwise-free-category-completion}
There is a canonical isomorphism of categories
\[
\begin{tikzcd}[column sep=large]
    \operatorname{Free}(G_L^\Sigma)
        \arrow[r, "\cong"]
    &
    \operatorname{Exec}(L).
\end{tikzcd}
\]
\end{theorem}

\begin{proof}
Define a functor
\[
    F:
    \operatorname{Free}(G_L^\Sigma)
    \to
    \operatorname{Exec}(L)
\]
which is the identity on residual languages and sends
\[
    (H;a_1,\ldots,a_n)
\]
to
\[
    (H,a_1\cdots a_n).
\]
The empty path is sent to \((H,\varepsilon)\), the identity execution at
\(H\).

For a fixed residual state \(H\), a letter sequence determines a unique path
in \(G_L^\Sigma\), because every successive vertex is forced to be
\[
    H,
    \partial_{a_1}H,
    \ldots,
    \partial_{a_1\cdots a_n}H.
\]
Every word in \(\Sigma^\ast\) has a unique decomposition as a finite letter
sequence. Hence \(F\) is bijective on every hom-set. Concatenation of letter
paths is sent to multiplication of words, so \(F\) preserves composition.
It is therefore an isomorphism.
\end{proof}

\begin{remark}[Letterwise paths and Segal completion]
\label{rem:letterwise-contraction}
The full word-labelled nerve is not a Segal composition structure on the fixed
letterwise trajectory sheaf. Indeed,
\[
    \mathscr T_L^\Sigma(1)
    =
    \operatorname{Der}(L)\times\Sigma,
\]
whereas
\[
    N_+\operatorname{Exec}(L)_1
    =
    \operatorname{Der}(L)\times\Sigma^\ast.
\]
A Segal composition structure on a fixed trajectory sheaf leaves the degree-one
object unchanged.

Rather, \(\operatorname{Exec}(L)\) is the free categorical completion of the
letterwise graph. It adjoins identity arrows labelled by the empty word and
composite one-step arrows labelled by arbitrary words. Its augmented nerve is
the canonical strict Segal object of that completed category, while
\[
    \mathscr T_L^\Sigma
    \hookrightarrow
    \operatorname{Path}(\operatorname{Exec}(L))
\]
is the subsheaf of paths segmented entirely into letters.
\end{remark}

\subsection{Classical finite-trajectory minimization}
\label{subsec:classical-temporal-minimization}

\begin{corollary}[Finite-trajectory minimization of deterministic automata]
\label{cor:classical-temporal-minimization}
Let \(\mathcal A\) be a reachable deterministic automaton recognizing \(L\).
Its semantic quotient onto the derivative automaton induces a surjection on
executions of every finite path length, both for the full
\(\Sigma^\ast\)-labelled execution semantics and for the letterwise
restriction. If \(\mathcal A\) is minimal, these maps are bijections.

For an arbitrary deterministic automaton recognizing \(L\), the same
statements hold after restriction to its reachable subautomaton.
\end{corollary}

\begin{proof}
For a reachable deterministic automaton recognizing \(L\), the semantic image
of the state map is precisely \(\operatorname{Der}(L)\). The semantic map is
therefore a surjective strict homomorphism onto the derivative automaton.
Theorem~\ref{thm:temporal-preservation-regular-quotients} gives surjections in
every finite path length for the full word-labelled execution semantics.

For letterwise paths, deterministic path lifting identifies a length-\(n\)
execution with an initial state together with an element of \(\Sigma^n\). The
induced map is the product of the state quotient with the identity on
\(\Sigma^n\), and is therefore surjective.

If the automaton is minimal, its semantic state map is an isomorphism, so the
induced maps on both full and letterwise executions are bijections. An
arbitrary recognizing automaton may contain unreachable states whose residual
behaviours are not derivatives of the language from the chosen initial state;
restriction to the reachable subautomaton removes this obstruction.
\end{proof}

\section{Relation to interval-based system semantics}
\label{sec:comparison-interval-sheaf-machines}

\subsection{Common span architecture}
\label{subsec:common-span-architecture-temporal}

Many sheaf-theoretic approaches to systems represent a machine by a span
\[
\begin{tikzcd}[column sep=large]
    \text{input trajectories}
    &
    \text{system trajectories}
        \arrow[l]
        \arrow[r]
    &
    \text{output trajectories}.
\end{tikzcd}
\]
The construction of this chapter has exactly that architecture. A Mealy
morphism determines
\[
\begin{tikzcd}[column sep=large]
    \operatorname{Path}(\mathbb A^{\mathrm{op}})
    &
    \mathscr T_\Phi
        \arrow[l, "{\operatorname{Path}(I_\Phi)}"']
        \arrow[r, "{\operatorname{Path}(O_\Phi)}"]
    &
    \operatorname{Path}(\mathbb B^{\mathrm{op}}).
\end{tikzcd}
\]
Composition is implemented by pullback over the shared intermediate
trajectory sheaf:
\[
\begin{tikzcd}[column sep=large]
    \mathscr T_{\Psi\circ\Phi}
        \arrow[r, "\cong"]
    &
    \mathscr T_\Phi
        \times_{\operatorname{Path}(\mathbb B^{\mathrm{op}})}
    \mathscr T_\Psi.
\end{tikzcd}
\]
This parallel is structural rather than merely notational. Both constructions
treat a system trajectory as a mediator between compatible input and output
trajectory data, and both use pullback to synchronize systems along a shared
interface.

\subsection{Derived rather than primitive temporal behaviour}
\label{subsec:difference-primitive-structure-temporal}

The difference lies in what is taken as primitive. In an interval-based sheaf
model, temporally extended behaviour is usually specified directly. One begins
with sections over intervals and restriction maps between them, and local
gluing is part of the system definition.

In the present theory, the primitive automata data are
\begin{enumerate}
    \item a persistent state carrier;
    \item typed input and output anchors;
    \item a deterministic transition operation;
    \item an output comparison satisfying the Mealy cocycle law.
\end{enumerate}
The trajectory sheaf is then derived from the execution category. Its gluing
property is not an independent axiom on the Mealy morphism: it follows from the
iterated pullback description of composable one-step executions.

This distinction has two consequences. First, every Mealy morphism yields a
discrete finite-time machine of interval-sheaf form. Second, an arbitrary
interval-based machine need not admit a presentation by a persistent state
carrier and deterministic one-step transition-output structure. Such a
presentation would require an additional reconstruction theorem, generally
with hypotheses expressing determinism, state persistence, and one-step
generation.

The present site has no durations or translations. A section over \([n]\)
records a path with \(n\) chosen categorical segments. Passing to interval
sites with lengths, translations, continuous durations, or hybrid time would
require an explicit change-of-site construction and a comparison of the
corresponding gluing conditions.

\subsection{The added categorical composition structure}
\label{subsec:segal-refinement-temporal}

The finite-path sheaf remembers temporal restriction and gluing. Because it is
derived from an internal category, it also has coherent categorical
composition data. The distinction can be summarized as follows:
\[
\begin{tikzcd}[column sep=large]
    N_+\operatorname{Exec}(\Phi)
        \arrow[r, "{j^\ast}"]
    &
    \mathscr T_\Phi.
\end{tikzcd}
\]
The object on the right records
\begin{enumerate}
    \item finite paths;
    \item restriction to contiguous temporal windows;
    \item initial and terminal truncation;
    \item gluing along matching state boundaries.
\end{enumerate}
The object on the left additionally records
\begin{enumerate}
    \item insertion of identity executions;
    \item contraction of adjacent execution steps;
    \item the unit and associativity equations for iterated contraction.
\end{enumerate}

This added structure is significant for the semantic applications.  Total
contraction is what reduces finite-path endpoint reachability to one-step
derivative generation.  In the later logical development, prefix contraction
will likewise compare universal finite-trajectory conditions with one-step
invariance.  Neither reduction would be available from a bare graph
trajectory sheaf.

\subsection{Scope and limitations}
\label{subsec:limitations-temporal-trajectories}

The constructions of this chapter are restricted to
\begin{enumerate}
    \item discrete linear path shapes;
    \item finite executions;
    \item deterministic Mealy transition structure;
    \item strict semantic homomorphisms;
    \item sequential composition of execution steps.
\end{enumerate}
They do not by themselves provide
\begin{enumerate}
    \item continuous-time or hybrid trajectory semantics;
    \item infinitary execution objects;
    \item branching, cubical, pomset, or higher-dimensional execution shapes;
    \item a classification of all interval-based machines by Mealy morphisms;
    \item a temporal fixed-point logic;
    \item a primitive notion of duration attached to a categorical arrow.
\end{enumerate}

The next chapter evaluates the output-comparison labels of these execution
categories in a downstream action.  The resulting relative execution category
will admit both an uncurried action-category presentation and a curried
polynomial total-response presentation.  The subsequent chapter will develop
the corresponding span, profile, and finite-trajectory logics.

\section{Chapter synthesis and passage to response semantics}
\label{sec:temporal-trajectories-chapter-synthesis}

\subsection{From automata to path sheaves}
\label{subsec:temporal-classification-summary}

The chapter began with the Mealy transition-output laws rather than with an
abstract interval site. From those laws we constructed the execution category
and the labelled span
\[
\begin{tikzcd}[column sep=large]
    \mathbb A^{\mathrm{op}}
    &
    \operatorname{Exec}(\Phi)
        \arrow[l, "I_\Phi"']
        \arrow[r, "O_\Phi"]
    &
    \mathbb B^{\mathrm{op}}.
\end{tikzcd}
\]
Its finite path objects are the finite runs of the automaton. Deterministic
path lifting identifies them as
\[
    \mathscr T_\Phi(n)
    \cong
    P\times_{A_0}
    \operatorname{Path}(\mathbb A^{\mathrm{op}})(n).
\]
Thus a run is uniquely determined by its initial state and input path.

Abstracting the restriction and gluing operations gave the finite-path site.
The resulting classification may be summarized as follows.

\begin{theorem}[Finite-path classification]
\label{thm:temporal-classification-summary}
The following hold:
\begin{enumerate}
    \item the augmented discrete interval category carries the topology
    generated by elementary-edge covers;
    \item an \(\mathcal E\)-valued presheaf is a trajectory sheaf exactly
    when
    \[
        X(-1)\cong1
    \]
    and, for every \(n\geq1\),
    \[
        X(n)
        \cong
        \underbrace{
        X(1)\times_{X(0)}\cdots\times_{X(0)}X(1)
        }_{n\text{ factors}};
    \]
    \item internal graphs are equivalent to trajectory sheaves,
    \[
    \begin{tikzcd}[column sep=large]
        \mathbf{Graph}_{\mathrm{int}}(\mathcal E)
            \arrow[r, "\simeq"]
        &
        \mathbf{TrajSh}_{\mathcal E};
    \end{tikzcd}
    \]
    \item adjacent path gluing is the canonical pullback isomorphism
    \[
        X(m+n)
        \cong
        X(m)\times_{X(0)}X(n).
    \]
\end{enumerate}
\end{theorem}

\begin{proof}
These are Proposition
\ref{prop:concrete-trajectory-sheaf-condition}, Corollary
\ref{cor:adjacent-temporal-gluing}, and Theorem
\ref{thm:graph-trajectory-sheaf-equivalence}.
\end{proof}

\subsection{From path sheaves to categorical execution}
\label{subsec:segal-execution-summary}

The graph-level path objects do not determine identity arrows or composition.
Those operations were recovered by extending from contiguous interval maps to
all order-preserving maps.

\begin{theorem}[Segal execution classification]
\label{thm:segal-execution-summary}
The following hold:
\begin{enumerate}
    \item augmented strict Segal objects are equivalent to internal categories,
    \[
    \begin{tikzcd}[column sep=large]
        \mathsf{Seg}^{\mathrm{str}}_+(\mathcal E)
            \arrow[r, "\simeq"]
        &
        \mathsf{IntCat}(\mathcal E);
    \end{tikzcd}
    \]
    \item the underlying trajectory sheaf of an internal category is the path
    sheaf of its underlying graph;
    \item internal category structures on a graph are equivalent to Segal
    composition structures on its trajectory sheaf;
    \item every Mealy morphism determines both a labelled trajectory span and
    a Segal-complete labelled execution span;
    \item path length records a chosen categorical segmentation, while the
    Segal contractions record its composite.
\end{enumerate}
\end{theorem}

\begin{proof}
These statements are Theorem
\ref{thm:internal-nerve-equivalence}, Proposition
\ref{prop:category-paths-underlying-graph}, Theorem
\ref{thm:classification-segal-enhancements}, and the constructions of
Section~\ref{sec:temporal-execution-semantics}.
\end{proof}

\subsection{Composition, minimization, and the next semantic layer}
\label{subsec:temporal-semantics-minimization-summary}

\begin{theorem}[Finite-trajectory semantics and minimization]
\label{thm:temporal-semantics-minimization-summary}
For Mealy morphisms and strict semantic homomorphisms:
\begin{enumerate}
    \item execution categories and trajectory sheaves are functorial in strict
    semantic homomorphisms;
    \item Mealy composition is represented by pullback,
    \[
        \mathscr T_{\Psi\circ\Phi}
        \cong
        \mathscr T_\Phi
        \times_{\operatorname{Path}(\mathbb B^{\mathrm{op}})}
        \mathscr T_\Psi;
    \]
    \item the composition comparison is natural and coherently associative;
    \item strict regular semantic quotients induce regular epimorphisms in
    every finite path length;
    \item behavioural minimization induces a canonical trajectory quotient
    \[
        \mathscr T_\Phi
        \twoheadrightarrow
        \mathscr T_{\operatorname{Beh}(\Phi)};
    \]
    \item reachable and separated realizations have the same trajectory sheaf
    as the canonical derivative realization;
    \item derivative-invariant predicates are precisely those closed under all
    finite residual prefixes;
    \item generated derivative closure is finite-path endpoint reachability.
\end{enumerate}
\end{theorem}

\begin{proof}
Functoriality and composition are Proposition
\ref{prop:functoriality-execution-categories}, Theorem
\ref{thm:execution-category-composition}, and Theorem
\ref{thm:temporal-composition}.  Quotient preservation and minimization are
Theorems \ref{thm:temporal-preservation-regular-quotients} and
\ref{thm:temporal-semantic-quotient}.  The residual statements are Theorems
\ref{thm:finite-horizon-characterization-contracts} and
\ref{thm:generated-closure-temporal-reachability}.
\end{proof}

\begin{remark}[Conceptual summary]
\label{rem:temporal-final-conceptual-summary}
The chapter has three interacting levels.

At the automata level, a Mealy morphism determines a category of one-step
executions and a labelled finite-run semantics.  At the graph level, finite
runs restrict to contiguous windows and glue uniquely from compatible
one-step segments.  At the categorical level, identity insertion and coherent
contraction organize those paths into an augmented strict Segal object.

The semantic consequence
\[
    \text{behavioural state minimization}
    \quad\Longrightarrow\quad
    \text{finite-trajectory minimization}
\]
is therefore compatible with all finite labelled execution structure.  For
full categorical execution, finite path length measures segmentation, not
primitive duration.  For letterwise Boolean execution, the generating graph
supplies a genuine atomic-step interpretation, and the full word-labelled
execution category is its free categorical completion.
\end{remark}

\subsection{Passage to evaluated and curried response semantics}
\label{subsec:temporal-passage-response-semantics}

The output of this chapter is the labelled execution span
\[
\begin{tikzcd}[column sep=large]
    \mathbb A^{\mathrm{op}}
    &
    \operatorname{Exec}(\Phi)
        \arrow[l, "I_\Phi"']
        \arrow[r, "O_\Phi"]
    &
    \mathbb B^{\mathrm{op}},
\end{tikzcd}
\]
together with its path sheaf and Segal composition structure.

The next chapter performs two operations on this data.  First, a right
\(\mathbb B\)-action evaluates the output-comparison label and produces a
relative execution category.  Second, the labelled one-step action is
transposed by dependent product into a polynomial total-response coalgebra.
The resulting comparison has the form
\[
\begin{tikzcd}[column sep=huge]
    \text{labelled relative action}
        \arrow[r, shift left=1.1ex, "\text{curry}"]
    &
    \text{multiplicative response coalgebra}
        \arrow[l, shift left=1.1ex, "\text{flatten}"']
\end{tikzcd}
\]
and flattening recovers the relative execution category.  The following logic
chapter will then compare predicates on selected execution witnesses,
predicate liftings on total response profiles, and modalities over the finite
trajectory objects constructed here.
